\documentclass[11pt]{article}
\usepackage[margin=1.05in]{geometry}
\usepackage{amsmath,amsthm,amssymb,mathtools,bm}
\usepackage{booktabs,array,enumitem}\usepackage[expansion=false]{microtype}
\usepackage{xcolor}
\usepackage[colorlinks=true,linkcolor=blue!60!black,citecolor=blue!60!black,urlcolor=blue!60!black]{hyperref}
\theoremstyle{plain}
\newtheorem{theorem}{Theorem}[section]\newtheorem{proposition}[theorem]{Proposition}\newtheorem{lemma}[theorem]{Lemma}
\newtheorem{corollary}[theorem]{Corollary}\newtheorem{conjecture}[theorem]{Conjecture}
\theoremstyle{definition}\newtheorem{definition}[theorem]{Definition}\newtheorem{remark}[theorem]{Remark}
\newtheorem{example}[theorem]{Example}
\newcommand{\R}{\mathbb R}\newcommand{\E}{\mathbb E}\newcommand{\Var}{\operatorname{Var}}\newcommand{\Cov}{\operatorname{Cov}}
\newcommand{\relu}{\operatorname{relu}}\newcommand{\diag}{\operatorname{diag}}\newcommand{\tr}{\operatorname{tr}}
\newcommand{\Tr}{\operatorname{Tr}}\newcommand{\kap}{\kappa}\newcommand{\Sym}{\operatorname{Sym}}
\newcommand{\He}{\mathrm{He}}\newcommand{\1}{\mathbf 1}\newcommand{\cA}{\mathcal A}\newcommand{\cB}{\mathcal B}
\newcommand{\cT}{\mathcal T}\newcommand{\cH}{\mathcal H}\newcommand{\cN}{\mathcal N}
\newcommand{\st}[1]{\textsc{#1}}
\title{\textsc{Localization in the commutant, stage 10:}\\ \large the selected Dirac operator.\\
Tilings, kink currents, contact memory and the He-initialised \textsc{ReLU} network\\ as a critical random spectral geometry}
\author{}
\date{9 October 2026}
\begin{document}
\maketitle
\begin{abstract}
We place the bias-free He-initialised \textsc{ReLU} multilayer perceptron with standard Gaussian input inside the
framework of random finite spectral triples (Dirac ensembles), Sahasrabudhe's two-level picture of probabilistic
combinatorics at small scales, and Klartag's evolving ellipsoid. The placement has one central object, the
\emph{selected Dirac operator}: the weights form a random chiral Dirac chain on the layered neuron set, the input
selects a projection of the commutative neuron algebra (its gate pattern), and the network is the resolvent of the
Dirac chain compressed by that projection. The partition of input space by the selected projection is the
activation tiling; its walls are an isotropic random hyperplane process at every layer (level one), and its overlap
kernels are the geometry in which the next layer's walls live (level two, the lift). Proved results:
(i) a Crofton formula for the tiling and the folding law at He criticality (lengths preserved, chords contract like
$3\pi/l$); (ii) the \emph{kink-current formula}: the exact mean of the network is a sum over walls of the Gaussian
surface density of the wall times the downstream transport, equivalently half the expected local time of the
pre-activations along a Brownian path; (iii) the third-cumulant memory of every moment chain is a record of
contacts: its hub index is the index of a kink and its identity leg the pinned coordinate (Price's theorem), and
a frame bound shows that its compressibility is that of its non-Hadamard leg; (iv) the neuron algebra and the
modular frame of the layer state are complementary, which is the Hadamard obstruction in operator-algebraic form;
(v) the Euler--Stein correction of stage 9 is an average over Gaussian balls of uniformly distributed radius in the
tiling, its constant is $1/(d+2)$ for a defect of Hermite contact order $d$, which is the order of the
concentration-function coefficient at the kink that the estimator gets wrong; (vi) neuron-pair correlations are
Mehler factors, i.e.\ imaginary modular times of the Gibbs state of the Hermite number operator, and the depth flow
of the pair temperature is marginal at linear order with a $\beta^{3/2}$ term produced by the kink, which fixes the
exponent of the folding law; (vii) the mean is the unique solution of a positivity-plus-Stein moment problem whose
first loop equation is the kink-current formula, and the cumulant chains are its truncations without positivity.
(viii) in the simplest framing (input slots as vertices, layers as quadratic forms, tokens as Hermite quanta) the
mean of every neuron is twice the trace of its quadratic form, layer two's form is the graph of two-step walks through
the hidden layer, cumulants are walk sums on the hidden codegree graph, the Gaussian closure is the two-token walk with
exact contacts, and the third-cumulant memory is first-order exchange whose hub coefficients are themselves walks.
Every statement is labelled as proved, known (with a reference), or conjectural; the identities are checked by
accompanying scripts.
\end{abstract}
\tableofcontents
\section{The placement in one page}
\label{sec:placement}
Fix $W_1,\dots,W_L\in\R^{n\times n}$ with i.i.d.\ $N(0,2/n)$ entries, put $h_0=x\sim\gamma=N(0,I_n)$, $z_l=W_lh_{l-1}$,
$h_l=\relu(z_l)$; the target is $u=\E h_L$. The placement consists of four identifications, each of which is a
statement proved below.

\paragraph{1. The space and the metric.} The layered neuron set $X=\{0,\dots,L\}\times[n]$ carries the commutative
algebra $\cA=C(X)$ and the Hilbert space $\ell^2(X)$. The weights are one self-adjoint operator
$D=\sum_l\big(W_l\otimes|l\rangle\langle l{-}1|+W_l^{\top}\otimes|l{-}1\rangle\langle l|\big)$, odd for the grading
$(-1)^l$: a random \emph{chiral Dirac chain}. He initialisation is the Gaussian (quadratic-action) Dirac ensemble
$e^{-\frac n8\Tr D^2}dD$ at the coupling where the gated transfer has unit mean-square gain (Proposition~\ref{prop:critical}).

\paragraph{2. The selection.} An input $x$ selects a projection $P_x\in\cA$ (its gate pattern), and the network is
the compressed resolvent $z_L(x)=\big[(1-D_{\downarrow}P_x)^{-1}\big]_{L0}\,x$ (Theorem~\ref{thm:resolvent}). The
\emph{selected Dirac operator} is $P_xDP_x$, the Dirac operator of the active subnetwork. Input space is tiled by
the value of the selected Dirac operator (the activation tiling), the Gaussian input induces a probability law on
compressions (a Dirac ensemble \emph{induced by the input}), and moving the input across a wall toggles one site,
a rank-one contact update of the selected geometry (Proposition~\ref{prop:contact}). Localization of the input
(stage 9) is the dynamical selection: a Gaussian ball of shrinking radius at a moving centre, whose terminal
selected Dirac operator is that of a Gaussian point.

\paragraph{3. The two levels.} At every layer the walls $\{z_{l,a}=0\}$ are the hyperplanes of $n$ i.i.d.\
isotropic normals in the previous feature space (level one: a random point process on a sphere). Their overlaps
---codegrees of caps, intersections of half-spaces---are arc-cosine kernels on both the input side and the neuron
side (Lemma~\ref{lem:duality}), and that level-two geometry is the ambient geometry of the next layer's walls (the
lift). Integral geometry (Crofton) counts the walls a curve crosses (Theorem~\ref{thm:crofton}); at He criticality
lengths are preserved and chords contract like $3\pi/l$ (Theorem~\ref{thm:folding}): the selected geometry folds.
The same two-level structure reappears one level up: kinks are points with marks (kink density, covariance arms),
and the third-cumulant memory is a sum of rank-one atoms over these points whose overlaps are codegrees in the
covariance graph (Section~\ref{sec:memory}).

\paragraph{4. The observable and its estimators.} The mean is a boundary functional of the tiling: the
\emph{kink current} $u=\sum_{l,a}\E\big[\delta(z_{l,a})\,|\nabla z_{l,a}|^2\,\partial_{h_l}h_L\,e_a\big]$
(Theorem~\ref{thm:kink}), the first Schwinger--Dyson equation of the network's Gaussian field theory
(Section~\ref{sec:bootstrap}). A moment chain is an \emph{unsharp} selection (gate projections replaced by their
expectations), its third cumulant is a record of contacts (hub $=$ kink, identity leg $=$ pinned coordinate,
Theorem~\ref{thm:birth}), and its errors are read by a clock: Gaussian balls of uniform radius in the tiling,
equivalently an exponentially distributed Ornstein--Uhlenbeck time, which turns Hermite contact order $d$ into the
Euler--Stein constant $1/(d+2)$ (Theorems~\ref{thm:uniform}, \ref{thm:order}).

\begin{center}\small
\begin{tabular}{>{\raggedright\arraybackslash}p{3.5cm}>{\raggedright\arraybackslash}p{3.5cm}>{\raggedright\arraybackslash}p{3.6cm}>{\raggedright\arraybackslash}p{4.0cm}}\toprule
Sahasrabudhe \cite{survey} & Klartag \cite{klartag} & Dirac ensembles \cite{BG,KP,DG,BDG} & He-\textsc{ReLU} network\\\midrule
Poisson points in a box & lattice points & random Dirac $D$ & weight rows: random walls per layer\\
conflict graph, codegree & contacts $\langle y,Ay\rangle=1$ & spectral triple $(\cA,\cH,D)$ & tiling; codegree $=$ arc-cosine kernel\\
lift (forget embedding) & ellipsoid $A_t$ (metric) & metric from $D$ (Connes) & selected Dirac $P_xDP_x$; pullback metric\\
nibble: $\Var\le p\sum c_y^2$ & rank-one contact updates & action $\Tr D^2,\Tr D^4$ & He $=$ quadratic action, critical gain\\
small codegree $\Rightarrow$ independence & volume growth & 1-cut/2-cut, symmetry breaking & folding law; odd memory sector\\
small-ball $\rho_h$ at scale $h$ & stopping when held & spectral dimension, variance & concentration coefficients at the kink\\
\bottomrule\end{tabular}
\end{center}
Two cautions recorded by the stage-9 workflow and kept here. The gate algebras are commutative subalgebras of
$L^\infty(\R^n,\gamma)$, so the modular theory of their states is trivial; the noncommutative geometry of the
problem lives in the pair (neuron algebra, Dirac/transport) and in the pair (neuron algebra, layer state on $M_n$),
Sections~\ref{sec:triple} and~\ref{sec:complement}. And Klartag's process is the dual of a localization: there the
metric is the martingale and the centre is fixed, here the centre is the martingale and the metric (precision)
deterministic; the correspondence is between the contact structures, not the processes.

\section{Two levels: random walls and their lift}
\label{sec:levels}
\subsection{Level one: the walls of a layer are an isotropic point process}
Condition on layers $1,\dots,l-1$. The feature map $x\mapsto h_{l-1}(x)\in\R^n$ is then fixed, and the rows
$w_{l,a}$ of $W_l$ are $n$ independent $N(0,\tfrac2nI)$ vectors; their directions $\hat w_{l,a}$ are $n$ i.i.d.\
uniform points on $S^{n-1}$, a binomial process that is Poisson to first order. Each point carries a \emph{wall} in
input space, $F_{l,a}=\{x:\langle w_{l,a},h_{l-1}(x)\rangle=0\}$, and each input $x$ carries a \emph{cap}, the open
hemisphere $C_l(x)=\{u\in S^{n-1}:\langle u,h_{l-1}(x)\rangle>0\}$; neuron $a$ is active at $x$ iff
$\hat w_{l,a}\in C_l(x)$. This is the incidence structure of \cite[\S3.3]{survey} (points in a box, balls around
them, intersections) with the r\^oles of the two families symmetric.

\begin{lemma}[Sign duality and the two codegrees]\label{lem:duality}
Write $\theta(u,v)$ for the angle between two vectors and $J(\theta)=\sin\theta+(\pi-\theta)\cos\theta$.
\begin{enumerate}[label=(\roman*),nosep]
\item (input side) For inputs $x,y$ with $\theta=\theta(h_{l-1}(x),h_{l-1}(y))$, conditionally on layers $<l$ the
codegree $N_l(x,y)=\#\{a:\text{$a$ active at $x$ and at $y$}\}$ is $\mathrm{Bin}(n,\tfrac{\pi-\theta}{2\pi})$, and the
number of neurons whose gate differs at $x$ and $y$ is $\mathrm{Bin}(n,\tfrac\theta\pi)$. The weighted codegree
$K_l(x,y)=\langle h_l(x),h_l(y)\rangle$ has conditional mean $\tfrac1\pi|h_{l-1}(x)||h_{l-1}(y)|\,J(\theta)$ and
conditional variance $O(1/n)$ relative to its mean squared.
\item (neuron side) At layer $1$, for two rows with angle $\theta_{ab}$, $\gamma(\text{$a$ and $b$ active})=
\tfrac{\pi-\theta_{ab}}{2\pi}$ and $\E_x[h_{1,a}h_{1,b}]=\tfrac1{2\pi}|w_a||w_b|\,J(\theta_{ab})$.
\item (Poisson covariance) If the number of points is Poissonised to $\mathrm{Poi}(n)$, the counts of points in two
caps have $\Cov(N(C_l(x)),N(C_l(y)))=n\,\mu(C_l(x)\cap C_l(y))=n\tfrac{\pi-\theta}{2\pi}$: the overlap of two caps is the
covariance of their counts, and the codegree is the pairing $\tau(p_xp_y)$ of the two cap projections in
$L^\infty(S^{n-1},\mu)$.
\end{enumerate}
\end{lemma}
\begin{proof}
For a uniform direction $u$ and two vectors at angle $\theta$, $\mathbb P(\langle u,v_1\rangle>0,\langle u,v_2\rangle>0)
=\tfrac{\pi-\theta}{2\pi}$ (the lune between two hemispheres) and $\mathbb P(\text{signs differ})=\tfrac\theta\pi$; the
rows are independent, which gives the binomial laws. The weighted codegree is a sum of $n$ i.i.d.\ terms
$\relu(\langle w,h(x)\rangle)\relu(\langle w,h(y)\rangle)$, whose mean is the degree-one arc-cosine kernel
$\tfrac2n\tfrac{|h(x)||h(y)|}{2\pi}J(\theta)$ \cite{CS}; the variance statement is the law of large numbers for $n$
i.i.d.\ terms. (ii) is the same computation with $x\sim\gamma$ isotropic and the rows fixed. (iii) is the Poisson
covariance identity: split the two caps into the three disjoint pieces; counts of disjoint pieces are independent and
the variance of a Poisson count equals its mean.
\end{proof}

So the post-activation covariance of a layer (the object every moment chain propagates) is the weighted codegree of
the neurons' half-spaces, and the feature kernel on inputs is the weighted codegree of the inputs' caps. Neither
side is privileged: the incidence $(x,a)\mapsto\1[\langle w_a,h(x)\rangle>0]$ is the sign pattern of one bilinear
form.

\subsection{Crofton: how many walls a curve crosses}
\begin{theorem}[Crofton formula for the tiling]\label{thm:crofton}
Let $c:[0,1]\to\R^n$ be a piecewise $C^1$ curve of inputs with $h_{l-1}(c(t))\neq0$, and let
$\Lambda_{l-1}(c)=\int_0^1\big|\tfrac{d}{dt}\,\widehat{h_{l-1}(c(t))}\big|\,dt$ be the angular length of its image
in the layer-$(l-1)$ feature sphere. Conditionally on layers $<l$, the expected number of crossings of layer-$l$ walls
by $c$ is
\[
\E\big[\#\{(a,t):z_{l,a}(c(t))=0\}\,\big|\,\text{layers}<l\big]\;=\;\frac{n}{\pi}\,\Lambda_{l-1}(c).
\]
\end{theorem}
\begin{proof}
For one row, $X_t=\langle w,\widehat{h_{l-1}(c(t))}\rangle\sqrt{n/2}$ is a centred Gaussian process with unit
variance and $\Cov(X_t,\dot X_t)=0$ (the derivative of a unit vector is orthogonal to it), so the Kac--Rice formula
gives $\E\#\{t:X_t=0\}=\tfrac1\pi\int_0^1\sqrt{\Var\dot X_t}\,dt=\tfrac1\pi\int_0^1|\tfrac{d}{dt}\widehat{h_{l-1}}|\,dt$.
The zeros of $X_t$ are the crossings of the wall $F_{l,a}$; sum over the $n$ rows.
\end{proof}

This is the integral-geometric face of the two-level picture: the expected number of random walls (level one)
that meet a set is the set's mean width in the geometry of the previous level. For a chord the count of
\emph{differing gates} is $n\theta/\pi$ (Lemma~\ref{lem:duality}); for the curve it is $n\Lambda/\pi$. The two differ as
soon as the image curve folds.

\subsection{The lift at He criticality: lengths are preserved, chords contract}
\begin{proposition}[Unit angular gain]\label{prop:gain}
Let $v\perp c$ be a tangent vector at a feature vector $c$, $c'=\relu(Wc)$, $\dot c'=D(Wc)\,Wv$ its image,
$W$ with i.i.d.\ $N(0,\sigma_w^2/n)$ entries. Then $\E|c'|^2=\tfrac{\sigma_w^2}2|c|^2$, $\E|\dot c'|^2=
\tfrac{\sigma_w^2}2|v|^2$, $\E\langle c',\dot c'\rangle=0$ and $\Var\langle c',\dot c'\rangle=O(|c|^2|v|^2/n)$. Hence the
expected squared angular speed is multiplied by $1+O(1/n)$ for every $\sigma_w$, and lengths are preserved
($\sigma_w^2=2$) iff the gated transfer has unit mean-square gain, $\chi_1=\sigma_w^2\,\E[\relu'(z)^2]=1$.
\end{proposition}
\begin{proof}
$\langle w_a,c\rangle$ and $\langle w_a,v\rangle$ are independent Gaussians because $v\perp c$; each neuron is active
with probability $\tfrac12$ independently of $\langle w_a,v\rangle$. Then $\E|\dot c'|^2=n\cdot\tfrac12\cdot
\tfrac{\sigma_w^2}n|v|^2$, $\E|c'|^2=n\cdot\E\relu(g)^2$ with $g\sim N(0,\sigma_w^2|c|^2/n)$, and
$\langle c',\dot c'\rangle=\sum_a\relu(\langle w_a,c\rangle)\langle w_a,v\rangle$ is a sum of $n$ independent centred terms
of variance $\tfrac{\sigma_w^2|c|^2}{2n}\tfrac{\sigma_w^2|v|^2}{n}$.
\end{proof}

\begin{theorem}[Folding law; infinite width]\label{thm:folding}
In the limit $n\to\infty$ with $L$ fixed (the NNGP limit \cite{CS,PLRSG}), the angle $\theta_l$ between the features
of two inputs obeys the $\sigma_w$-independent recursion $\cos\theta_{l+1}=J(\theta_l)/\pi$, and
\[
\theta_{l+1}=\theta_l-\frac{\theta_l^2}{3\pi}+O(\theta_l^3),\qquad
\theta_l=\frac{3\pi}{\,l+3\pi/\theta_0+O(\log l)}\;.
\]
The angular length of every curve is preserved, $\Lambda_l=\Lambda_0$. Consequently, for two inputs at angle
$\theta_0$, the walls of layer $l$ crossed by the great-circle arc between them number $n\theta_0/\pi$ in
expectation at every layer (in total $L n\theta_0/\pi$), while the gates that differ between the endpoints number
$n\theta_{l-1}/\pi\approx3n/l$ (in total $\approx3n\log L$). The selected geometry is an isometric immersion of the
input sphere into the feature sphere whose chords shrink like $3\pi/l$: it folds.
\end{theorem}
\begin{proof}
The recursion is Lemma~\ref{lem:duality}(i) with the law of large numbers, normalised by the norms (the $\sigma_w$
factor cancels in the angle because $\relu$ is positively homogeneous). Expanding,
$\cos\theta'-\cos\theta=(\sin\theta-\theta\cos\theta)/\pi=\theta^3/(3\pi)+O(\theta^5)$, and
$\cos\theta'=1-\theta'^2/2+O(\theta'^4)$ give $\theta'^2=\theta^2-\tfrac{2\theta^3}{3\pi}+O(\theta^4)$, i.e.\
$\theta'=\theta-\theta^2/(3\pi)+O(\theta^3)$. Then $1/\theta_{l+1}=1/\theta_l+1/(3\pi)+O(\theta_l)$, and summing the
$O(\theta_l)=O(1/l)$ errors gives the $O(\log l)$ term. The angle map $f(\theta)=\arccos(J(\theta)/\pi)$ has
$f'(0)=1$, so the length of a curve, $\lim\sum f(\Delta\theta_i)=\sum\Delta\theta_i$, is preserved; this is the
infinite-width form of Proposition~\ref{prop:gain}. The crossing counts follow from Theorem~\ref{thm:crofton} and
Lemma~\ref{lem:duality}(i).
\end{proof}

\begin{remark}[Two metrics on inputs]
The \emph{chord} metric $|\hat h_l(x)-\hat h_l(y)|=2\sin(\theta_l/2)$ and the \emph{length} metric (the infimum of
image lengths of curves joining $x$ and $y$) separate at depth: the first collapses like $3\pi/l$, the second is
the round metric of the input sphere at every depth. Section~\ref{sec:distance} identifies the second as the
Connes distance of the selected Dirac operator. Bias-free \textsc{ReLU} networks are thus at the edge of chaos for
every $\sigma_w$ in the angle variable, and He initialisation is the point at which the length variable is critical too
\cite{PLRSG,HDR}; the folding exponent is computed spectrally in Proposition~\ref{prop:heattrace}. The linear-in-$L$
crossing count is the infinite-width form of the region counts of \cite{HR}.
\end{remark}

\subsection{The lift as a tower of algebras}
Let $\cB_l\subset L^\infty(\R^n,\gamma)$ be the von Neumann algebra generated by the gate functions
$\1[z_{k,a}>0]$, $k\le l$. Then $\cB_1\subset\cB_2\subset\cdots\subset\cB_L$; $\cB_l$ is the algebra of functions
constant on the tiles of the depth-$l$ tiling, its minimal projections are the tiles, and the trace $\gamma$ gives
each tile its Gaussian measure. The level-two data of layer $l$ (the kernels $K_l$ on inputs, the post-activation
covariance $C_l$ on neurons) are the ambient metric data of the level-one point process of layer $l+1$. The network
is a $\cB_L$-valued piecewise-linear section: on the tile with projection $p_g$ it is the linear map $J_g$, so that
$h_L=\sum_gp_g\,J_gx$. Everything that follows is about how the walls of this tower, and not its tiles, carry the
mean.

\section{The network as a random spectral triple}
\label{sec:triple}
\subsection{The chiral Dirac chain}
\begin{definition}[Network spectral triple]\label{def:triple}
Let $X=\{0,\dots,L\}\times[n]$, $\cA=C(X)\cong\bigoplus_{l=0}^L\mathbb C^n$ acting diagonally on $\cH=\ell^2(X)$, and
\[
D=D_\downarrow+D_\downarrow^{\top},\qquad D_\downarrow=\sum_{l=1}^LW_l\otimes|l\rangle\langle l-1|,
\]
with grading $\Gamma=\bigoplus_l(-1)^l$ and real structure $\mathcal J=$ complex conjugation. Then $D\Gamma=-\Gamma D$,
$\mathcal JD=D\mathcal J$, $\mathcal J\Gamma=\Gamma\mathcal J$, $\mathcal J^2=1$, and the order-zero and first-order
conditions hold trivially because $\cA$ is commutative: $(\cA,\cH,D,\Gamma,\mathcal J)$ is a real even finite spectral
triple of KO-dimension $0$. Its commutators $[D,f]$, $f\in\cA$, have entries $W_l(a,b)\,(f_{l-1,b}-f_{l,a})$: $D$ is
the weighted adjacency of the layered graph, and the Connes distance of the triple is the metric on neurons in
which strong weights are short edges.
\end{definition}

\begin{proposition}[He initialisation as a critical Gaussian Dirac ensemble]\label{prop:critical}
If the entries of every $W_l$ are i.i.d.\ $N(0,\sigma_w^2/n)$, the law of $D$ on the space of chiral chain operators is
$Z^{-1}\exp\!\big(-\tfrac{n}{4\sigma_w^2}\Tr D^2\big)\,dD$, the quadratic Barrett--Glaser action \cite{BG,KP} restricted to
this kinematic space. For a fixed vector $z$ and a gate projection $P$ selecting the positive coordinates of $z$,
$\E|W\,Pz|^2=\sigma_w^2|Pz|^2$ and $\E|Pz|^2=\tfrac12|z|^2$ for symmetric $z$; the gated transfer has unit
mean-square gain exactly at $\sigma_w^2=2$, and $\E\log(|z_{l+1}|^2/|z_l|^2)=\log(\sigma_w^2/2)+O(1/n)$. He
initialisation is the critical coupling of the selected (gated) chain.
\end{proposition}
\begin{proof}
$\Tr D^2=2\sum_l\Tr W_l^{\top}W_l$ and the Gaussian density of the entries is $\prod\exp(-\tfrac{n}{2\sigma_w^2}W_{ij}^2)$.
For fixed $y=Pz$, $Wy$ has i.i.d.\ $N(0,\sigma_w^2|y|^2/n)$ coordinates. The logarithmic statement is the law of
$\log\chi^2_n/n$, whose mean is $-O(1/n)$.
\end{proof}

\subsection{The network is a compressed resolvent}
\begin{theorem}[Gated resolvent]\label{thm:resolvent}
For an input $x$ let $P_x\in\cA$ be the projection that is the identity on layer $0$ and the gate projection
$\1_{(0,\infty)}(z_l(x))$ on layer $l\ge1$, and $\psi=\bigoplus_lz_l$ the field of pre-activations ($z_0=x$). Then $\psi$ is
the unique solution of the self-consistent equation
\[
\psi=x\oplus0+D_\downarrow P_\psi\,\psi,\qquad\text{i.e.}\qquad \psi=(1-D_\downarrow P_\psi)^{-1}(x\oplus0),
\]
where $P_\psi=\1_{(0,\infty)}(\psi)$ is a spectral projection of the field in the commutative algebra. Since
$D_\downarrow P$ is nilpotent, $(1-D_\downarrow P)^{-1}=\sum_{k=0}^{L}(D_\downarrow P)^k$, and
\begin{gather*}
z_L(x)=\big[(1-D_\downarrow P_x)^{-1}\big]_{L0}\,x=W_LP_{L-1}W_{L-1}\cdots P_1W_1\,x,\\
\big[z_L(x)\big]_k=\sum_{a_0\to a_1\to\cdots\to a_L=k}\ \prod_{l=1}^LW_l(a_l,a_{l-1})\prod_{l=1}^{L-1}\1[z_{l,a_l}(x)>0]\;x_{a_0}.
\end{gather*}
\end{theorem}
\begin{proof}
The equation is block lower triangular: its layer-$l$ component reads $\psi_l=W_lP_{l-1}\psi_{l-1}$ with
$P_{l-1}\psi_{l-1}=\relu(\psi_{l-1})$, which is the forward pass; uniqueness follows by induction on $l$. The Neumann
series terminates because $(D_\downarrow P)^{L+1}=0$, and the $(L,0)$ block of $(D_\downarrow P)^L$ is the displayed
product, whose entries are the path sums.
\end{proof}

The network therefore computes with one fixed random operator $D$ and one input-dependent projection $P_x$ of the
commutative algebra. Everything nonlinear is in the selection $x\mapsto P_x$; given $P_x$ the network is the
resolvent of a compressed Dirac operator. The path expansion is the sum over walks of the counting problems of
stage 9's reading (exterior-algebra path counts): the gates are input-dependent vertex weights.

\subsection{The selected Dirac operator and the induced ensemble}
\begin{definition}[Selection]\label{def:selection}
The \emph{selected Dirac operator} of $x$ is $D_x=P_xDP_x$, the Dirac operator of the active subnetwork, and the
\emph{selected metric} is the pullback $g_x=J(x)^{\top}J(x)$ of the output metric, $J(x)=\partial h_L/\partial x$.
The activation tiling is the partition of $\R^n$ by the value of $P_x$; the input law induces on the finite set of
compressions $\{P_gDP_g\}$ the probability $\gamma(\text{tile }g)$. We call this law the Dirac ensemble \emph{induced
by the input}: a random finite geometry drawn from the $2^{n(L-1)}$ compressions of one random Dirac operator.
\end{definition}

\begin{proposition}[Walls are rank-one contacts]\label{prop:contact}
Let $x$ cross the wall $F_{l,a}$ transversally from $z_{l,a}<0$ to $z_{l,a}>0$, all other gates fixed. Then $P_x$
changes by the minimal projection $e_{(l,a)}$, the selected Dirac operator changes by
$e_{(l,a)}DP+PDe_{(l,a)}+e_{(l,a)}De_{(l,a)}$ (a perturbation of rank at most two supported on the edges at the toggled
site), the Jacobian changes by the rank-one operator
\[
\Delta J=\big(\partial_{h_l}h_L\big)e_a\,\nabla z_{l,a}^{\top},
\]
and the selected metric changes by $\Delta g=\nabla z_{l,a}\,v^{\top}J+J^{\top}v\,\nabla z_{l,a}^{\top}+|v|^2\nabla z_{l,a}
\nabla z_{l,a}^{\top}$ with $v=(\partial_{h_l}h_L)e_a$, a rank-two update whose new direction is the normal of the wall.
\end{proposition}
\begin{proof}
Only the gate of $(l,a)$ changes, so $J=\partial_{h_l}h_L\,P_l\,\partial_xz_l$ changes by $\partial_{h_l}h_L\,e_ae_a^{\top}
\partial_xz_l$, and $e_a^{\top}\partial_xz_l=\nabla z_{l,a}^{\top}$; downstream gates are continuous across a generic
wall. The other two statements are direct expansions.
\end{proof}

This is the precise sense in which the tiling plays the r\^ole of Klartag's lattice. In \cite{klartag} the ellipsoid
matrix $A_t$ moves until a lattice point $y$ reaches its boundary, and the contact imposes the rank-one constraint
$\Tr(A\,yy^{\top})=1$. Here the selected geometry is constant inside a tile and receives a rank-one (Jacobian) or
rank-two (metric) update at every wall a path crosses; along any path $x_t$ the selected operators $D_{x_t}$ form a
Hamiltonian on a dynamically evolving graph whose edges switch at wall crossings. The localization processes of stage
9 are exactly such paths with a Gaussian ball attached: on Eldan's path the ball $N(m_t,(1+t)^{-1}I)$ shrinks around a
martingale centre $m_t$ until it lies in one tile, and the terminal selected operator is $D_X$, $X\sim\gamma$. The
process is dual to Klartag's (there the metric is the martingale and the centre fixed); what transfers is the contact
structure, and Section~\ref{sec:clock} shows that the estimator errors are read off exactly from it.

\subsection{Distance on states}\label{sec:distance}
\begin{proposition}[Connes--Kantorovich distance of the selected geometry]\label{prop:distance}
Let
\[d_J(x,y)=\inf_c\int_0^1|J(c(t))\,\dot c(t)|\,dt\]
over piecewise $C^1$ curves $c$ from $x$ to $y$ (the length of the image curve; a pseudometric because $J$ may have a
kernel), and for probability laws $\nu,\nu'$ with first moments let
\[
d_D(\nu,\nu')=\sup\{|\nu(f)-\nu'(f)|:\ f\ \text{$1$-Lipschitz for }d_J\}.
\]
Then $d_D=W_1^{d_J}$ (Kantorovich--Rubinstein), every output coordinate is $1$-Lipschitz for $d_J$, and hence
$|u(\nu)_k-u(\nu')_k|\le W_1^{d_J}(\nu,\nu')$ for the exact means of any two input laws. For a smooth nondegenerate
metric $d_D$ is Connes' spectral distance with the Dirac operator of $g$ \cite{Connes,Rief}; here it is its
length-space extension to the piecewise-flat selected metric. In the infinite-width limit $d_J$ restricted to the
input sphere is the round metric at every depth (Theorem~\ref{thm:folding}), while the chord metric of the features
collapses: the distance on states is preserved by the selection and the distance between features is not.
\end{proposition}
\begin{proof}
Kantorovich--Rubinstein duality on the (pseudo)metric space $(\R^n,d_J)$; $|h_{L,k}(x)-h_{L,k}(y)|\le|h_L(x)-h_L(y)|\le
\int|J\dot c|$ for every curve. The last statement is Theorem~\ref{thm:folding}.
\end{proof}

\section{The mean is a kink current}
\label{sec:current}
\subsection{The kink-current formula}
Write $\partial_{h_l}h_L(x)$ for the Jacobian of the output with respect to the post-activations of layer $l$
($=I$ for $l=L$, and $D_LW_LD_{L-1}\cdots W_{l+1}$ otherwise, with gates at $x$), and $f_{l,a}$ for the density of
$z_{l,a}(x)$, $x\sim\gamma$.

Call a tile \emph{dead} if some hidden layer is entirely inactive on it ($h_l\equiv0$ there); on a dead tile the network
vanishes identically and several walls coincide along its faces. The Gaussian measure of the dead tiles is
$\mathbb P(\exists\,l<L:\ h_l(x)=0)$, an orthant probability that is exponentially small in $n$ ($1.6\cdot10^{-1}$ at
$n=6$, $8\cdot10^{-6}$ at $n=16$, $L=3$; of order $2^{-n}$ in general and below $10^{-300}$ at $n=1024$).

\begin{theorem}[Kink current]\label{thm:kink}
For almost every weight realisation, up to a remainder carried by the faces of the dead tiles,
\begin{equation}\label{eq:kink}
\begin{split}
u=\E\,h_L(x)&=\sum_{l=1}^{L}\sum_{a=1}^{n}\E\Big[\delta\big(z_{l,a}(x)\big)\,\big|\nabla z_{l,a}(x)\big|^2\,
\partial_{h_l}h_L(x)\,e_a\Big]\\
&=\sum_{l,a}f_{l,a}(0)\;\E\Big[|\nabla z_{l,a}|^2\,\partial_{h_l}h_L\,e_a\ \Big|\ z_{l,a}=0\Big].
\end{split}
\end{equation}
Each term is the Gaussian surface integral of the wall $F_{l,a}$, weighted by the squared gradient of the
pre-activation that defines it and by the transport of a unit kick at $(l,a)$ to the output.
\end{theorem}
\begin{proof}
$F=h_L$ is piecewise linear and positively homogeneous, so Euler's relation $F=x\cdot\nabla F$ holds off the walls, and
$\nabla F=J^{\top}$ is piecewise constant with jumps across the walls. Gaussian integration by parts for the bounded
variation field $J$ gives $\E F=\E[x\cdot\nabla F]=\E[\operatorname{div}\nabla F]=\E[\Delta F]$, where $\Delta F$ is the
distributional Laplacian, a measure carried by the walls. Write $J=D_LW_LD_{L-1}\cdots D_1W_1$ and differentiate the
gates: $\partial_iD_l=\sum_a\delta(z_{l,a})\,\partial_iz_{l,a}\,e_ae_a^{\top}$. By the product rule
\[
\textstyle\sum_i(\partial_iJ)_{\cdot i}=\sum_{l,a}\delta(z_{l,a})\,\big(\partial_{h_l}h_L\big)e_a\;\sum_i\partial_iz_{l,a}\,
\big(e_a^{\top}W_lD_{l-1}\cdots W_1\big)_i=\sum_{l,a}\delta(z_{l,a})\,|\nabla z_{l,a}|^2\,\big(\partial_{h_l}h_L\big)e_a ,
\]
because $e_a^{\top}W_lD_{l-1}\cdots W_1=\nabla z_{l,a}^{\top}$. Distinct walls meet in sets of codimension two for almost
every weight realisation, the downstream gates are continuous across a generic point of $F_{l,a}$, and
$\{z_{l,a}=0\}$ is a finite union of pieces of hyperplanes (on each live tile $z_{l,a}$ is a nonzero linear form
almost surely), so the products of distributions are well defined and $\E[\delta(z)G]=f(0)\,\E[G\mid z=0]$ with a continuous
density $f$. Taking expectations gives \eqref{eq:kink}. The only exception is a face of a dead tile: there a
downstream pre-activation vanishes identically on one side, several gates jump at once, the product rule must be
applied to the simultaneous jump, and $z_{l,a}$ has an atom at $0$; these faces carry the remainder, which is bounded by the
Gaussian surface measure of the dead tiles' boundaries times the size of the transport. (The identity check confirms
both halves: at $n=16$ the two sides agree to the Monte Carlo noise, $1.5\%$; at $n=6$, where dead tiles carry $16\%$
of the mass, they differ by $25\%$.)
\end{proof}

\begin{example}
For $L=1$: $f_{1,a}(0)=1/(\sqrt{2\pi}|w_a|)$, $|\nabla z_{1,a}|^2=|w_a|^2$, so $u_a=|w_a|/\sqrt{2\pi}$, the mean of
$\relu(N(0,|w_a|^2))$. For $L=2$ the layer-$1$ terms give $D_2W_2$ applied to the layer-$1$ kink masses on the walls, and
the layer-$2$ terms give the kink masses of the curved walls $\{w_{2,a}\cdot h_1(x)=0\}$.
\end{example}

\begin{theorem}[Tanaka form: the mean is transported local time]\label{thm:tanaka}
Let $B$ be a standard Brownian motion in $\R^n$ started at $0$ and $L^{(l,a)}_t$ the semimartingale local time at $0$ of
$z_{l,a}(B_t)$. Then $h_L(B_t)$ is a semimartingale with
\begin{gather*}
dh_L(B_t)=J(B_t)\,dB_t+\tfrac12\sum_{l,a}\big(\partial_{h_l}h_L\big)(B_t)\,e_a\,dL^{(l,a)}_t,\\
u=\E h_L(B_1)=\tfrac12\sum_{l,a}\E\int_0^1\big(\partial_{h_l}h_L\big)(B_t)\,e_a\,dL^{(l,a)}_t .
\end{gather*}
In particular the \emph{kink mass} $m_{l,a}=\E[\delta(z_{l,a})|\nabla z_{l,a}|^2]$ is half the expected local time that
$z_{l,a}$ accumulates at its kink along a unit-time Brownian path from the origin.
\end{theorem}
\begin{proof}
$h_L$ is a difference of convex functions on each line, piecewise linear, and the It\^o--Tanaka (Meyer) formula
applies: $dh_L(B)=\nabla h_L\cdot dB+\tfrac12\Delta h_L(dt)$ with the Laplacian measure computed in the proof of
Theorem~\ref{thm:kink}; the occupation-time formula converts $\delta(z_{l,a})|\nabla z_{l,a}|^2dt$ into $dL^{(l,a)}$.
Homogeneity gives $\E[\delta(z(B_t))|\nabla z(B_t)|^2G(B_t)]=t^{-1/2}\E[\delta(z)|\nabla z|^2G]$ for degree-zero $G$, and
$\tfrac12\int_0^1t^{-1/2}dt=1$, consistent with \eqref{eq:kink}.
\end{proof}

\begin{corollary}[Mean is created only at kinks]\label{cor:created}
For every neuron, $\E h_{l,a}=m_{l,a}+\E\big[\1[z_{l,a}>0]\,\Delta z_{l,a}\big]$, where $\Delta z_{l,a}$ is the
wall measure of the lower layers transported to $z_{l,a}$. A layer adds mean exactly through its kink masses and
transports the lower layers' kink currents through its gates; a network without kinks has zero mean
($\E z_l=\E\Delta z_l$ and $\Delta$ of a linear map vanishes).
\end{corollary}
\begin{proof}
Apply the computation of Theorem~\ref{thm:kink} to the single coordinate $h_{l,a}$:
$\Delta h_{l,a}=\delta(z_{l,a})|\nabla z_{l,a}|^2+\1[z_{l,a}>0]\Delta z_{l,a}$.
\end{proof}

\subsection{What an estimator has to compute: conditional expectations on walls}
Formula \eqref{eq:kink} reduces the mean to two kinds of quantities per wall: the density $f_{l,a}(0)$ of the
pre-activation at its kink, and a conditional expectation \emph{on the wall}, a measure-zero event. The first is the
$h\to0$ limit of the concentration function $Q_{l,a}(h)=\mathbb P(|z_{l,a}|\le h)=2f_{l,a}(0)h+\tfrac13f''_{l,a}(0)h^3+\dots$
(the pre-activation law is absolutely continuous with a continuous density, small-ball exponent one, no arithmetic
structure: the gates are functions of a continuous Gaussian). On the official network $0$ the Gaussian surrogate
$\varphi(\alpha)/\sigma$ of $f_{l,a}(0)$ is right to $1$--$3\cdot10^{-3}$ at every layer when the mean and variance are
right (stage-9 workflow, $2.6\cdot10^5$ inputs), so the densities are not where estimators fail. The second kind is
a conditional expectation relative to an event of probability zero, the extreme case of the rare-event-relative
conditioning of \cite[\S4]{survey} and \cite{SST}: the error tolerated in $\E[\,\cdot\mid z_{l,a}=0]$ is relative to
the kink mass, not absolute. For a Gaussian state, conditioning on $z_{l,a}=0$ is a rank-one \emph{pinning}
(a directional localization along $\nabla z_{l,a}$ taken to its end), and Section~\ref{sec:memory} shows that the
third cumulant of every moment chain is the record of exactly these pinnings.

\section{Estimators are unsharp selections; their memory is a record of contacts}
\label{sec:memory}
\subsection{Moment chains are unsharp selections}
A moment chain does not know $P_x$; it knows the Gaussian state $N(m_l,S_l)$ of each layer and replaces the gate
projection by its expectation, the diagonal effect $\Phi_l=\diag\,\Phi(\alpha_{l,a})\in\cA$, $0\le\Phi_l\le1$. Its
transport is $T_{l\leftarrow l'}=W_l\Phi_{l-1}W_{l-1}\cdots\Phi_{l'}W_{l'+1}$, the resolvent of the Dirac chain compressed by an
effect instead of a projection.

\begin{proposition}[Unsharp selection]\label{prop:unsharp}
With the path notation of Theorem~\ref{thm:resolvent},
\[
\E_x\big[\partial_{h_{l'}}z_l(x)\big]-T_{l\leftarrow l'}=\sum_{\text{paths}}\ \prod W\ \Big(\E\big[\textstyle\prod_k\1[z_{k,a_k}>0]\big]-\prod_k\Phi_{k,a_k}\Big),
\]
the sum over paths of the connected gate correlations along the path. The chain is exact for the expected
transport iff the gates on every path are uncorrelated; its correction terms are correlations between the
projection and the field that selects it (measurement back-action).
\end{proposition}
\begin{proof}
Take expectations in the path expansion term by term.
\end{proof}

Only the on-wall transports of \eqref{eq:kink} enter the mean, so what matters is the back-action \emph{conditional on a
kink}. That is where the third cumulant comes from.

\subsection{Pinning and births: the hub is a kink}
\begin{lemma}[Pinning]\label{lem:pinning}
If $z\sim N(m,S)$, then $z$ conditioned on $z_a=0$ is $N\big(m-S e_a\,m_a/S_{aa},\;S-Se_ae_a^{\top}S/S_{aa}\big)$: a
rank-one update of the state along the covariance column of the pinned neuron. Downstream gate expectations shift
to first order by $-\varphi(\alpha_b)\,\rho_{ab}\,\alpha_a$, $\rho_{ab}=S_{ab}/\sqrt{S_{aa}S_{bb}}$.
\end{lemma}
\begin{proof}
Gaussian conditioning; $\Phi\big((m_b-S_{ab}m_a/S_{aa})/\sigma_b\big)-\Phi(\alpha_b)=-\varphi(\alpha_b)\rho_{ab}\alpha_a+O(\rho^2)$.
\end{proof}

\begin{theorem}[Births are contacts]\label{thm:birth}
Let $z\sim N(m,S)$ with small off-diagonal covariances and $h=\relu(z)$. For distinct $i,j,k$,
\[
\kap_3(h_i,h_j,h_k)=\sum_{c\in\{i,j,k\}}\ \Phi(\alpha_a)\,S_{ac}\ \frac{\varphi(\alpha_c)}{\sigma_c}\ S_{cb}\,\Phi(\alpha_b)\Big|_{\{a,b\}=\{i,j,k\}\setminus\{c\}}+O(S_{\rm off}^3),
\]
the sum over which of the three neurons is the \emph{hub}. The hub contributes through $\E[\relu''(z_c)]=\varphi(\alpha_c)/\sigma_c=
f_c(0)$, its kink density, and the two arms are the covariances to the hub weighted by the gates. In CP form,
$\kap_3=\Sym\sum_c\,(\Phi\circ S_{\cdot c})\otimes 3e_c\otimes(f_c(0)\,S_{c\cdot}\circ\Phi)$ to second order in $S_{\rm off}$, which is the
star block of the organisers' chain $(\Phi\circ S^{\rm off},\,3I,\,(w_2\circ S^{\rm off}\circ\Phi)^{\top})$, $w_2=f(0)$; its
residual block makes the birth exact beyond second order.
\end{theorem}
\begin{proof}
Joint cumulants of functions of a Gaussian vector expand over connected graphs on the index set with edges weighted by
off-diagonal covariances and vertex factors $\E f^{(\deg)}(z_v)$ (Price's theorem, or Wick's theorem applied to the
Hermite expansions). On three vertices the connected graphs with two edges are the three paths, centred at the hub,
with vertex factors $\E\relu'(z_a)=\Phi(\alpha_a)$, $\E\relu''(z_c)=\varphi(\alpha_c)/\sigma_c$, $\E\relu'(z_b)=\Phi(\alpha_b)$; the
triangle is $O(S_{\rm off}^3)$. Writing the hub as the middle leg $e_c$ gives the CP form, and $\Sym$ distributes the three
positions (factor $3$).
\end{proof}

So the identity leg that makes the third-cumulant memory expensive is not an accident of a factorisation: it is the
label of the kink at which the cumulant was born. Every birth block is $n$ rank-one atoms, one per neuron $c$ of
the birth layer, with intensity $f_c(0)$ and arms $S_{\cdot c}$; transported, the atom of hub $c$ born at $l'$ and read at
$l$ is $u_c z_c^{\top}$ with $u_c=(T A e_c)\circ(T e_c)$ and $z_c=TCe_c$, $T=T_{l\leftarrow l'}$. Lemma~\ref{lem:pinning}
says what the atom is: the first-order response of the downstream state to pinning neuron $c$ at its kink.
Measured on network $0$ (stage-9 workflow): the atom energies are proportional to $f_c(0)^2$ (log-log correlation
$0.97$--$0.998$ at ages $1,3,6$), and the atoms are mutually orthogonal to $10^{-3}$
($\|\sum_cP^{(c)}\|^2/\sum_c\|P^{(c)}\|^2=0.998$--$1.004$).

\subsection{Kinks as points: codegree and the frame bound}
The second appearance of the two levels: the kinks of a layer are a marked point process on the neuron index
(marks: density $f_c(0)$, arms $S_{\cdot c}$), and the memory is the sum of their atoms. Two kinks overlap through the
neurons their arms share, $\langle S_{\cdot c}\circ\Phi,S_{\cdot c'}\circ\Phi\rangle=(S\Phi^2S)_{cc'}$: a codegree in the
covariance graph. At He initialisation the off-diagonal covariances are $O(n^{-1/2})$, so codegrees are $O(n^{-1/2})$ of
the degrees, and the Gram matrix of the atoms, a Hadamard product of two such codegree matrices, has off-diagonal
entries $O(n^{-1})$ of the diagonal. This is exactly the small-codegree regime of \cite[\S3.3--3.4]{survey}. There it
is what makes the nibble work (the conditional variance of a surviving degree is at most $p\sum_yc_y^2$); here the
same orthogonality makes the memory incompressible.

\begin{theorem}[Frame bound]\label{thm:frame}
Let $P=UZ^{\top}=\sum_{c=1}^Nu_cz_c^{\top}$ with $U,Z\in\R^{n\times N}$, and let $\hat P=UMZ^{\top}$ with $\operatorname{rank}M\le k$ be
any compression in the hub index (keeping $k$ hubs, reweighting, or mixing hubs through any $k$-dimensional subspace).
Then
\[
\|P-\hat P\|_F^2\ \ge\ \lambda_{\min}(U^{\top}U)\,\|(I-M)Z^{\top}\|_F^2\ \ge\ \lambda_{\min}(U^{\top}U)\,\lambda_{\min}(Z^{\top}Z)\,(N-k).
\]
\end{theorem}
\begin{proof}
$\|U(I-M)Z^{\top}\|_F^2=\tr\big(Z(I-M)^{\top}G_u(I-M)Z^{\top}\big)\ge\lambda_{\min}(G_u)\|(I-M)Z^{\top}\|_F^2$ with $G_u=U^{\top}U$, and
$\|(I-M)Z^{\top}\|_F^2=\tr\big((I-M)G_z(I-M)^{\top}\big)\ge\lambda_{\min}(G_z)\|I-M\|_F^2\ge\lambda_{\min}(G_z)(N-k)$ by
Eckart--Young.
\end{proof}

\begin{corollary}[The memory is as compressible as its transport leg, no more]
If the Hadamard leg is well conditioned, $U^{\top}U\approx cI$ (orthogonal hub atoms of comparable size), every
hub-index compression loses at least $c$ times the energy of the transport leg $Z=TC$ outside the kept subspace. The
cheapest compression is therefore the low-rank truncation of $TC$, and its error is that truncation's error. On network $0$
the Hadamard bulk is near-isotropic while $TC$ keeps $91/96/99\%$ of the slice energy at rank $64/128/256$ (stage-9
workflow), against the $1-10^{-4}$ of energy that a $10^{-2}$ relative slice accuracy requires; this is the measured
failure of every tier in stage 9, now with its reason.
\end{corollary}

\subsection{Complementarity of the neuron algebra and the layer state}\label{sec:complement}
The gate algebras are commutative and the modular theory of their states is trivial. The genuinely noncommutative
pair is the neuron algebra $A=\diag(\R^n)\subset M_n$, where the nonlinearity acts (gates, Hadamard products,
functional calculus), and the layer state $\rho_l=M_l/\Tr M_l$ on $M_n$, $M_l=\E[z_lz_l^{\top}]$, faithful and
non-tracial, whose modular group $\sigma_s(Y)=\rho_l^{is}Y\rho_l^{-is}$ is the frame in which the transport acts.

\begin{proposition}[Takesaki for the neuron algebra]\label{prop:takesaki}
The only conditional expectation of $M_n$ onto $A$ is the pinching $\Delta(Y)=\diag(Y)$, and it preserves $\rho$ iff
$\rho$ is diagonal; equivalently (Takesaki \cite{Tak}) $A$ is $\sigma^\rho$-invariant iff $\rho\in A$. The independent-neuron
(mean-field) coarse-graining is state-preserving only for uncorrelated layers.
\end{proposition}
\begin{proof}
For a conditional expectation $E$ onto $A$ and $i\neq j$, bimodularity gives $E(e_{ij})=e_{ii}E(e_{ij})e_{jj}=0$ since
$E(e_{ij})$ is diagonal; so $E=\Delta$. $\Tr(\rho\,\Delta Y)=\Tr(\Delta(\rho)Y)$, which equals $\Tr(\rho Y)$ for all $Y$ iff
$\rho=\Delta(\rho)$.
\end{proof}

\begin{theorem}[Complementarity]\label{thm:complement}
Let $B$ be the masa of the eigenprojections of $\rho_l$ (simple spectrum), $U$ its eigenbasis, and
$C_{ai}=|U_{ai}|^2-\tfrac1n$. With $\tau=\tfrac1n\Tr$, $E_A=\Delta$ and $E_B$ the pinching in the eigenbasis,
$\|E_BE_A-J\|_{L^2(\tau)\to L^2(\tau)}=\|C\|$, $J(Y)=\tau(Y)1$. Every $\rho_l$-preserving conditional expectation $E_\cN$
onto a subalgebra $\cN\subseteq B$ (any spectral coarse-graining of the layer state) factors through $E_B$, and for
every $Y\in A$ with $\tau(Y)=0$,
\[
\|E_B\,Y\|_2\le\|C\|\,\|Y\|_2 .
\]
If the eigenvectors are Haar-distributed on the orthogonal group, $\|C\|\approx2\sqrt{2/n}$.
\end{theorem}
\begin{proof}
The first identity is the alternating-expectation theorem for two partitions of unity into minimal projections
(verified in the stage-9 workflow; it follows from writing $E_A=VV^*$, $E_B=WW^*$ with isometries onto the partitions
with $V^*W=C+\tfrac1n\mathbf 1\mathbf 1^{\top}$, so that $E_BE_A-J=WC^{\top}V^*$ has norm $\|C\|$). Subalgebras of $B$ are invariant under
$\sigma^{\rho}$ because $\rho\in B$, so $\rho$-preserving expectations onto them exist; each is the identity on $B$-valued
data after $E_B$ (for $\cN$ spanned by blocks $Q_k$ of eigenprojections, $E^\rho_\cN(Y)=\sum_k\Tr(\rho Q_kY)/\Tr(\rho Q_k)\,Q_k$ and
$\Tr(\rho Q_kY)=\Tr(\rho Q_kE_B(Y))$). For $Y\in A$ with $\tau(Y)=0$, $E_BY=(E_BE_A-J)Y$. For Haar orthogonal eigenvectors $n|U_{ai}|^2$ is asymptotically $\chi^2_1$, so $nC$ has
centred entries of variance $2$ and the norm of an $n\times n$ such matrix is $\approx2\sqrt{2n}$.
\end{proof}

At He initialisation the neuron algebra and the modular frame of the layer state are nearly mutually unbiased: a
compression that respects the state (keeps eigen-directions of the covariance) sees a neuron-diagonal observable
(a gate, a Hadamard product, a slice's hub diagonal) only to $O(n^{-1/2})$ ($2\sqrt{2/n}\approx0.09$ at $n=1024$), and a compression that respects the
neurons does not preserve the state. This is the Hadamard obstruction of stage 9 in operator-algebraic form, and it
is why the spectral birth truncations of stage 9 saturate (rank $16$ already keeps the spectral content and still
loses a factor three): spectral information and readout information live in complementary masas.

\subsection{Parity: the odd sector and broken symmetry}
\begin{proposition}[Odd sector]\label{prop:odd}
Fix layers $1,\dots,l-1$. The readout at layer $l$ (before its gates) of any third-cumulant atom born at $l'<l$ is a
homogeneous cubic polynomial in $W_l$, hence odd under $W_l\mapsto-W_l$, and its conditional mean over the Gaussian law
of $W_l$ vanishes. Any surrogate that depends on $W_l$ only through even statistics (spectra, $W_l^{\top}W_l$, the NNGP
recursion, free-probability models) assigns it zero.
\end{proposition}
\begin{proof}
Each of the three legs $TAe_c$, $Te_c$, $TCe_c$ contains $W_l$ exactly once as the last factor, and no gate follows it.
\end{proof}

In the language of \cite{DG}: the Gaussian weight ensemble has the symmetry $W_l\mapsto-W_l$ for every $l$, the
ensemble-averaged (free, NNGP) description is the symmetric phase, and each realisation carries an odd component of
self-averaging size and realisation-specific sign, as the broken-symmetry equilibrium measures of the $(1,0)$
ensemble do while the mean density of states stays symmetric. The analogy is exact in its consequence (ensemble
surrogates cannot see the odd sector; the stage-9 workflow measured that its effect on the final mean is still half
even, through odd readout coefficients, so surrogates capture at most a factor two of a memory worth a factor one
hundred) and inexact in its mechanism: here the breaking is explicit at relative order $n^{-1/2}$ (a random field),
not spontaneous at order one. The collective scalars that the chain must track exactly---the radial fourth-cumulant
scalar, worth a factor $16$, the trace and the collective eigenvalue of the layer covariance---play the r\^ole of the
multi-trace couplings $(\tr H^2)^2$, $\tr H\,\tr H^3$ of the Barrett--Glaser actions, whose values decide the phase there.

\section{The clock of the selection: Mehler factors, modular time and the Euler--Stein constant}
\label{sec:clock}
\subsection{The Euler--Stein correction probes the estimator on balls of uniform radius}
Stage 9 (Corollary \emph{Eldan's path}) gives, for a chain $E$ exact on point masses and scale covariant, with
$e(y)=E(N(y,I))$ and $\delta(y)=e(y)-y\cdot\nabla e(y)-\Delta e(y)$,
$u-e(0)=-\tfrac12\int_0^\infty(1+\tau)^{-3/2}\E_{y\sim N(0,\tau I)}\delta(y)\,d\tau$.

\begin{theorem}[Uniform radius]\label{thm:uniform}
With $r=(1+\tau)^{-1/2}\in(0,1]$,
\[
u-e(0)\;=\;-\int_0^1\E\Big[\delta\big(M_r/r\big)\Big]\,dr,\qquad M_r\sim N\big(0,(1-r^2)I\big).
\]
The localized laws on Eldan's path are the Gaussian balls $N(M_r,r^2I)$ of radius $r$ about a Gaussian centre, and under
the defect weight their radius is \emph{uniformly distributed} on $[0,1]$; equivalently the Ornstein--Uhlenbeck time
$t=-\log r$ of the localization is a standard exponential variable.
\end{theorem}
\begin{proof}
$d\tau=-2r^{-3}dr$, so $\tfrac12(1+\tau)^{-3/2}d\tau=\tfrac12r^3\cdot2r^{-3}|dr|=dr$. On Eldan's path the localized law at
time $\tau$ is $N(m,sI)$ with $s=(1+\tau)^{-1}=r^2$ and $y=m/\sqrt s\sim N(0,\tau I)$, i.e.\ $m\sim N(0,r^2\tau I)=N(0,(1-r^2)I)$.
Finally $\mathbb P(-\log r>t)=\mathbb P(r<e^{-t})=e^{-t}$.
\end{proof}

This is the precise form of ``an ellipsoid in a tiling'': the correction is the average, over Gaussian balls of
every radius $r\in(0,1)$ with equal weight, of the estimator's defect on that ball. A ball of radius $r$ meets a wall
through its centre's neighbourhood with probability proportional to $r$ (Crofton), and the order at which the
estimator mis-accounts the contact decides how fast its defect vanishes as the ball shrinks.

\subsection{Hermite contact order and the Euler--Stein constant}
\begin{theorem}[Contact order law]\label{thm:order}
Suppose the defect has the kink profile of order $d$ at a wall through the origin, $\delta(y)=c\,(\varphi\He_d)(\langle v,y\rangle)$
with $|v|=1$ (the form of every layer-one wall of a bias-free network). Then
$\E_{y\sim N(0,\tau I)}\delta(y)=r^{d+1}\delta(0)$, and the merge constant defined by $u-e(0)=-a\,\delta(0)$ is
\[
a=\int_0^1r^{d+1}\,dr=\frac1{d+2}=\E\big[e^{-(d+1)t}\big],\quad t\sim\mathrm{Exp}(1).
\]
For odd $d$ the defect vanishes at the origin and contributes nothing.
\end{theorem}
\begin{proof}
$\varphi\He_d=(-1)^d\varphi^{(d)}$, so $\E_{s\sim N(0,\tau)}(\varphi\He_d)(s)=(-1)^d(\varphi*\varphi_\tau)^{(d)}(0)=(-1)^d\varphi_{1+\tau}^{(d)}(0)$,
and $\varphi_{1+\tau}(x)=r\varphi(rx)$ gives $\varphi^{(d)}_{1+\tau}(0)=r^{d+1}\varphi^{(d)}(0)=(-1)^dr^{d+1}(\varphi\He_d)(0)$.
Integrate against $dr$ (Theorem~\ref{thm:uniform}).
\end{proof}

\begin{proposition}[Contact order is a concentration coefficient]\label{prop:conc}
For a pre-activation with Gaussian surrogate $N(m,\sigma^2)$, $\alpha=m/\sigma$, the concentration function at the kink is
\begin{gather*}
Q(h)=\mathbb P(|z|\le h)=2f(0)\,h+\tfrac13f''(0)\,h^3+O(h^5),\\
f(0)=\frac{\varphi(\alpha)}\sigma=\frac{(\varphi\He_0)(\alpha)}{\sigma},\qquad
f''(0)=\frac{(\varphi\He_2)(\alpha)}{\sigma^3}.
\end{gather*}
The coefficient of $h^{d+1}$ is the order-$d$ Hermite kink content. An estimator whose leading defect is an error in
that coefficient has, by Theorem~\ref{thm:order}, merge constant $1/(d+2)$: the Euler--Stein constant reads off which
Taylor coefficient of the concentration functions at the kinks the estimator gets wrong.
\end{proposition}
\begin{proof}
$Q(h)=\int_{-h}^hf=2f(0)h+\tfrac13f''(0)h^3+O(h^5)$, and $\partial_z^k\big[\varphi((z-m)/\sigma)/\sigma\big]_{z=0}=
(-1)^k\sigma^{-k-1}\varphi^{(k)}(-\alpha)$, with $\varphi^{(k)}=(-1)^k\He_k\varphi$ and parity.
\end{proof}

\begin{conjecture}[Contact order at depth]\label{conj:order}
For the chains of stage 9 the defect at depth is dominated by one contact order, and $a=1/(d+2)$ with: $d=0$ for the
Gaussian closure and for chains with truncated third-cumulant memory (their errors are variance errors acting through
$f(0)$), $d=2$ for the exact first-order third-cumulant chain (its errors are second-order, the curvature $f''(0)$).
\end{conjecture}
\begin{center}\small
\begin{tabular}{lccc}\toprule
chain (official networks, $2049$ exact chain runs each) & measured $a$ & predicted & order $d$\\\midrule
Gaussian closure, network $0$ / $1$ & $0.535$ / $0.500$ & $1/2$ & $0$\\
window-one third-cumulant chain (memory dropped beyond one age) & $0.482$ & $1/2$ & $0$\\
exact first-order third-cumulant chain, network $0$ / $1$ & $0.255$ / $0.257$ & $1/4$ & $2$\\\bottomrule
\end{tabular}
\end{center}
The layer-wise profile of the absolute kink density along the localization was measured at all sixteen layers of
network $0$ to decay as $(1+\tau)^{-1/2}=r$ (stage-9 workflow), which is the $d=0$ profile; no profile of the exact
chain's defect has been measured, so the $d=2$ attribution rests on the constant. The earlier explanation of
$0.255$ as ``two births, $1/3$, reduced by rotation'' is superseded: the order law gives $1/4$ with no correction.

\subsection{Mehler factors are imaginary modular times}
\begin{proposition}[Mehler--KMS]\label{prop:mehler}
Let $N$ be the Hermite number operator on $L^2(\R,\gamma)$ ($N\He_k=k\He_k$), $\omega_\beta(Y)=\Tr(\rho_\beta Y)$ with
$\rho_\beta=(1-e^{-\beta})e^{-\beta N}$ the Gibbs state of $N$, faithful and non-tracial. For a standard Gaussian pair
$(X,Y)$ with correlation $\rho=e^{-\beta}\in(0,1)$ and real $f,g\in L^2(\gamma)$,
\[
\E\,f(X)g(Y)=\langle f,e^{-\beta N}g\rangle=\frac1{1-\rho}\,\omega_\beta\big(|g\rangle\langle f|\big).
\]
The modular group of $\omega_\beta$ is $\sigma_s(Y)=e^{-i\beta sN}Ye^{i\beta sN}$, the Mehler rotation (fractional Fourier
transform of Gaussian space); the Ornstein--Uhlenbeck semigroup $e^{-tN}$ is its continuation to imaginary time. The
correlation of two Gaussian variables is thus the Boltzmann factor of one Hermite quantum at inverse temperature
$\beta=-\log\rho$, and every pair correlation of the network, every Euler--Stein weight and every Wick coefficient of
the chains is a matrix element of $e^{-\beta N}$ at some $\beta$.
\end{proposition}
\begin{proof}
Mehler's formula: $\E f(X)g(Y)=\sum_k\rho^k\langle f,H_k\rangle\langle g,H_k\rangle$ with $H_k=\He_k/\sqrt{k!}$; and
$\Tr(\rho_\beta|g\rangle\langle f|)=(1-e^{-\beta})\langle f,e^{-\beta N}g\rangle$. The modular group of a Gibbs state $\rho_\beta$ is
$\mathrm{Ad}\,\rho_\beta^{is}$.
\end{proof}

\begin{proposition}[The ReLU heat trace and the folding exponent]\label{prop:heattrace}
Let $K(\beta)=\E[\relu(X)\relu(Y)]$ at correlation $e^{-\beta}$, a weighted heat trace of the Hermite spectrum of
$\relu$: $K(\beta)=\sum_kc_k^2e^{-\beta k}/k!$ with $c_0=\tfrac1{\sqrt{2\pi}}$, $c_1=\tfrac12$, $c_{2j+2}=(-1)^j(2j-1)!!/\sqrt{2\pi}$,
$c_{\rm odd\ge3}=0$. The weights decay like $k^{-5/2}$ (the kink), and
\[
K(\beta)=\tfrac12\Big(1-\beta+\tfrac{2\sqrt2}{3\pi}\beta^{3/2}+O(\beta^2)\Big),\qquad
\beta_{l+1}=\beta_l-\tfrac{2\sqrt2}{3\pi}\beta_l^{3/2}+O(\beta_l^2)
\]
for the pair temperature $\beta_l=-\log\cos\theta_l$ of two inputs. The linear coefficient is exactly the critical one
(the depth flow is marginal at linear order because $\relu$ is homogeneous: $K'(0)/K(0)=-\E\relu'^2/\E\relu^2=-1$), and the
kink's $\beta^{3/2}$ term drives $\beta_l\simeq9\pi^2/(2l^2)$, i.e.\ $\theta_l\simeq3\pi/l$. Depth heats the neuron pairs
quadratically, $T_l=1/\beta_l\simeq2l^2/(9\pi^2)$.
\end{proposition}
\begin{proof}
$\E[\relu(Z)\He_k(Z)]=\E[\relu^{(k)}(Z)]=\E[\delta^{(k-2)}(Z)]=\He_{k-2}(0)\varphi(0)$ for $k\ge2$, and
$\He_{2j}(0)=(-1)^j(2j-1)!!$; Stirling gives $c_{2j+2}^2/(2j+2)!\sim(8\pi^{3/2})^{-1}j^{-5/2}$. In closed form
$K=\tfrac1{2\pi}(\sin\theta+(\pi-\theta)\cos\theta)$ with $\cos\theta=e^{-\beta}$, whose non-analytic part
$\tfrac1{2\pi}(\sin\theta-\theta\cos\theta)=\tfrac{\theta^3}{6\pi}+O(\theta^5)$ with $\theta=\sqrt{2\beta}(1+O(\beta))$ is
$\tfrac{\sqrt2}{3\pi}\beta^{3/2}$. The recursion is $e^{-\beta_{l+1}}=K(\beta_l)/K(0)$, and $\beta^{-1/2}$ grows by
$\tfrac{\sqrt2}{3\pi}$ per layer.
\end{proof}

\begin{remark}[One clock]
Three flows are one structure. Depth raises the pair temperature, $\beta_l\to0$ like $l^{-2}$, with exponent fixed by the
Weyl-type tail $k^{-5/2}$ of the activation's Hermite spectrum (a smooth activation has an analytic heat trace and flows
exponentially unless tuned). Localization lowers it: the ball radius is the Boltzmann factor $r=e^{-t}$ of the
Ornstein--Uhlenbeck time, and the Euler--Stein constant $1/(d+2)=\E e^{-(d+1)t}$ is the mean Boltzmann factor of $d+1$
quanta at an exponential time. The quantum Hermite transform of stage 9's reading is the real-time version of the same
group. In this sense the Gaussian input's modular flow, the network's depth flow and the estimator's error clock are
the imaginary-time, discrete-time and averaged-time faces of the single semigroup $e^{-tN}$.
\end{remark}

\section{Bootstrap: the mean as the unique solution of a positivity problem}
\label{sec:bootstrap}
Khalkhali and Pagliaroli \cite{KP} solve Dirac ensembles without solving them: Schwinger--Dyson (loop) equations give
linear relations among moments, positivity of moment matrices cuts out a convex island, and the island shrinks onto
the answer as the moment degree grows. The network admits the same treatment, with the Gaussian input as the
``action'', Stein's identity as the loop equation, and the gates as idempotent variables.

\begin{definition}[The ReLU moment problem]\label{def:moment}
Let $\R[x,\theta]$ be the polynomial algebra in $x\in\R^n$ and commuting idempotents $\theta_{l,a}$ ($\theta^2=\theta$),
$l=1,\dots,L$. Define recursively the polynomials $z_1=W_1x$, $z_{l,a}=\sum_bW_l(a,b)\,\theta_{l-1,b}\,z_{l-1,b}$. A
\emph{feasible moment functional} is a linear $\mathcal L:\R[x,\theta]\to\R$ with $\mathcal L(1)=1$, Gaussian
$x$-moments (equivalently the loop equations $\mathcal L(x_iq)=\mathcal L(\partial_iq)$ for $q\in\R[x]$), positive moment
matrix $\mathcal L(p^2)\ge0$, and positive localizing forms $\mathcal L(p^2\,\theta_{l,a}z_{l,a})\ge0$,
$\mathcal L(-p^2(1-\theta_{l,a})z_{l,a})\ge0$ for all $p$.
\end{definition}

\begin{theorem}[The island is a point]\label{thm:island}
For almost every weight realisation every feasible moment functional represented by a measure is the law of
$(x,\theta(x))$ except for the values of the gates on dead tiles, where all downstream pre-activations vanish and the
gates are unconstrained; these measures are moment-determinate, and they all give the same value to every
pre-activation moment. In particular $u=\mathcal L(\theta_L\circ z_L)$ is determined by positivity and the loop
equations alone.
\end{theorem}
\begin{proof}
A representing measure $\mu$ lives on $\R^n\times\{0,1\}^{nL}$ and is supported in
$K=\{\theta_{l,a}z_{l,a}\ge0,\ (1-\theta_{l,a})z_{l,a}\le0\}$. If $z_{1,a}(x)>0$ the second constraint forces
$\theta_{1,a}=1$, if $z_{1,a}(x)<0$ the first forces $\theta_{1,a}=0$; given the gates of layers $<l$, $z_l(x,\theta)$ is the
true pre-activation and the same argument forces layer $l$. So the fibre of $K$ over $x$ is the single point $\theta(x)$
unless some $z_{l,a}(x)=0$. On a live tile this happens on a $\gamma$-null set ($z_{l,a}$ is a nonzero linear form there
almost surely); on a dead tile, where $h_l\equiv0$, all later pre-activations vanish whatever the later gates are, so
the fibre is a set of gate patterns on which every $z_{l'}$, $l'>l$, and $h_L$ take the value $0$. Since the $x$-marginal
is $\gamma$, $\mu$ is the graph measure up to the choice of these irrelevant gates. Determinacy: $\mu=\sum_\vartheta\mu_\vartheta\otimes\delta_\vartheta$
over the finitely many gate patterns, with $\mu_\vartheta\le\gamma$; the moments of $\mu$ determine those of every $\mu_\vartheta$
(the monomials in $\theta$ span the indicators of patterns), and each $\mu_\vartheta$ satisfies Carleman's condition because its
moments are bounded by the Gaussian ones.
\end{proof}

Under the standard convergence theory of moment--sum-of-squares relaxations for determinate measures
\cite{Las}, the islands of degree $2k$ form a nested sequence of intervals containing $u$ and shrinking to it. Two
features are specific to the network.

\begin{proposition}[Loop equations need the kink generators]\label{prop:loops}
Stein's identity for test functions that involve gates produces the distributions $\delta(z_{l,a})$:
$\partial_i\theta_{l,a}=\delta(z_{l,a})\,\partial_iz_{l,a}$. The loop equations therefore close only on the algebra generated by
$x$, the gates and the kink distributions $\delta^{(k)}(z_{l,a})$. The first of them, for the test field $q=J_{\cdot i}$, is
the kink-current formula \eqref{eq:kink}; the Wick coefficients $\E[\partial^k\relu^p]$ of every cumulant chain are the
Gaussian-state values of these kink generators.
\end{proposition}
\begin{proof}
Differentiate $\theta_{l,a}(x)=\1[z_{l,a}(x)>0]$ in the sense of distributions; the first loop equation is
$\E[x\cdot\nabla h_L]=\E[\Delta h_L]$, computed in Theorem~\ref{thm:kink}. For a Gaussian $z$, $\E[\delta^{(k)}(z)]$ are the
derivatives of the density at the kink, which is what the chains' Wick tables contain.
\end{proof}

\begin{proposition}[The level-one island of a single neuron is Scarf's interval]\label{prop:scarf}
If only the mean $m$ and variance $s^2$ of $z_{l,a}$ are imposed, together with $h\ge0$, $h\ge z$, $h(h-z)=0$, the closure of the set of feasible
values of $\E h$ is exactly $\big[\max(0,m),\ \tfrac12(m+\sqrt{m^2+s^2})\big]$. The Gaussian closure picks the interior
point $s\,(\alpha\Phi(\alpha)+\varphi(\alpha))$, $\alpha=m/s$.
\end{proposition}
\begin{proof}
The constraints force $h=z_+$; the lower bound is Jensen, and the upper bound is Scarf's worst case over laws with
given mean and variance, attained by a two-point law \cite{Scarf}.
\end{proof}

So a moment chain is a truncated Schwinger--Dyson hierarchy solved with a Gaussian (or Edgeworth) ansatz and
\emph{without} positivity. Positivity is what would turn a chain into a certificate; its cost at width $1024$ is
prohibitive at any useful degree, and its role here is structural: it shows that the exact answer is the unique point
of an island cut out by the input's loop equations and the complementarity of the gates, and that the information a
chain discards is the positivity of the kink generators' moment matrices, which is exactly what goes wrong when a
closure produces a negative fourth cumulant or a variance outside its admissible range.

\section{The token picture: vertices, quadratic forms, two-token walks}
\label{sec:tokens}
This section develops the simplest framing of the problem: every input slot is a vertex, the weights are weighted
edges, a layer is a quadratic form on the vertex graph, the next layer is read from the two-token (Kikuchi) lift of
that graph, and a second tier---the input law---sets the environment in which the lower tier's walks take place. The
framing is exact: it is the Wiener chaos (bosonic Fock) structure of Gaussian space, with tokens as Hermite quanta.
It reorganises everything above in a few lines, and it locates the obstruction precisely: it is a contact of tokens at
a hidden neuron.

\subsection{Dictionary}
Let $V=[n]$ be the input slots. By the Wiener--It\^o isomorphism
\[
L^2(\R^n,\gamma)\;\cong\;\bigoplus_{k\ge0}\Sym^k\big(\ell^2(V)\big),\qquad
F=\sum_k\big\langle f_k,\,{:}x^{\otimes k}{:}\big\rangle,\qquad f_k=\tfrac1{k!}\,\E\big[\nabla^kF\big]\ \ \text{(Stroock)},
\]
the $k$-th chaos is spanned by the multivariate Hermite polynomials $\He_\alpha$, $|\alpha|=k$: configurations of $k$
indistinguishable tokens on the vertices, repetitions allowed (the bosonic $k$-token Kikuchi level). The Gaussian input
is the vacuum $\Omega=1$; $a_i=\partial_i$ removes a token at $i$, $a_i^\dagger=x_i-\partial_i$ adds one, $N=\sum_ia_i^\dagger a_i$ is the
Ornstein--Uhlenbeck number operator. A weighted graph on $V$ is a function on the pair groupoid $V\times V$, i.e.\ an element
of $M_n=\mathbb C[V\times V]$; it moves one token along an edge as $d\Gamma(A)$ and all tokens as $\Gamma(A)$. Two-token kernels
are quadratic forms, i.e.\ elements of the groupoid algebra, and contracting two kernels in $r$ shared token slots,
$f\otimes_rg$, is groupoid convolution ($Q\otimes_1Q'=QQ'$, $Q\otimes_2Q'=\Tr QQ'$). Positive definite functions on the pair
groupoid act by Schur multiplication as completely positive maps; the unit space (the diagonal) is where two tokens sit
on one vertex.

\subsection{Layer one: each neuron is a one-mode token tower, and its mean is a trace}
\begin{proposition}[Ridge towers]\label{prop:ridge}
With $c_k=\E[\relu(Z)\He_k(Z)]$, i.e.\ $(c_0,c_1,c_2,c_3,c_4,c_5,c_6,\dots)=(\varphi_0,\tfrac12,\varphi_0,0,-\varphi_0,0,3\varphi_0,\dots)$,
$\varphi_0=1/\sqrt{2\pi}$, $c_k=\He_{k-2}(0)\varphi_0$ for $k\ge2$,
\[
h_{1,a}=|w_a|\sum_k\frac{c_k}{k!}\big(a^\dagger(\hat w_a)\big)^k\Omega ,\qquad
f_2(h_{1,a})=Q_a=\frac{w_aw_a^{\top}}{2\sqrt{2\pi}\,|w_a|}.
\]
Each layer-one neuron lives in the one-mode Fock space of its own edge vector, and its quadratic form is the rank-one
graph of that edge. Its degree-$k$ coefficient is $\E[\relu^{(k)}]$: the gate for $k=1$, the kink density for $k=2$, and
derivatives of the kink for $k\ge3$.
\end{proposition}
\begin{proof}
$\relu(|w|s)=|w|\relu(s)$ with $s=\hat w\cdot x\sim N(0,1)$, the one-dimensional Hermite expansion, and $\He_k(\hat w\cdot x)=
(a^\dagger(\hat w))^k\Omega$. $c_k=\E[\relu^{(k)}]=\E[\delta^{(k-2)}]=\He_{k-2}(0)\varphi_0$ by Gaussian integration by parts.
\end{proof}

\begin{theorem}[Euler--Fock ladder]\label{thm:ladder}
If $F$ is positively homogeneous of degree one then $(N+\sum_ia_ia_i)F=F$, i.e.
\[
f_0=2\tr f_2,\qquad \tr_{23}f_3=0,\qquad f_k=-\frac{(k+1)(k+2)}{k-1}\,\tr_{k+1,k+2}f_{k+2}\quad(k\ge2).
\]
In particular the mean of every neuron of every layer is twice the trace of its quadratic form: the mean is a closed
two-token loop, the two tokens of the quadratic form placed on the same input vertex and summed.
\end{theorem}
\begin{proof}
$x\cdot\nabla=\sum_i(a_i^\dagger+a_i)a_i=N+\sum_ia_ia_i$, and Euler's relation $x\cdot\nabla F=F$ holds for one-homogeneous $F$.
$\sum_ia_ia_i$ maps the chaos-$(k+2)$ component with kernel $f_{k+2}$ to chaos $k$ with kernel $(k+2)(k+1)\tr f_{k+2}$, the trace over
two token slots. Compare the chaos-$k$ components. For $h_{1,a}$: $2\tr Q_a=|w_a|/\sqrt{2\pi}$, and $-12\tr f_4=f_2$ with
$f_4=-\tfrac{\varphi_0}{24}|w_a|\hat w_a^{\otimes4}$.
\end{proof}

\subsection{Layer two: the quadratic form is the graph of two-step walks}
\begin{theorem}[Layer-two tokens]\label{thm:layer2}
For $z_{2,c}=\sum_aW_2(c,a)h_{1,a}$, with $\omega=W_2(c,\cdot)^{\top}$,
\begin{gather*}
z_{2,c}=m_c+\langle f_c,x\rangle+{:}x^{\top}Q_cx{:}+(\text{chaos}\ge4),\\
f_c=\tfrac12W_1^{\top}\omega,\qquad
Q_c=W_1^{\top}D_cW_1,\qquad D_c=\diag\Big(\frac{W_2(c,a)}{2\sqrt{2\pi}|w_a|}\Big).
\end{gather*}
$Q_c(i,j)=\sum_aW_1(a,i)\,D_c(a)\,W_1(a,j)$ is the weighted graph of two-step walks $i\to a\to j$ through the hidden layer,
turning at $a$ with weight $D_c(a)$. By Sylvester's identity every spectral quantity of $Q_c$ is one of $D_cG$, where
$G=W_1W_1^{\top}$ is the codegree graph of the hidden neurons (overlaps of their input supports; normalised, the cosines
$\rho_{ab}$).
\end{theorem}
\begin{proof}
Linearity of the chaos decomposition and Proposition~\ref{prop:ridge} ($\E g_{1,a}=\tfrac12$).
\end{proof}

\begin{theorem}[Cumulants are walk sums]\label{thm:walks}
(i) For $z=m+\langle f,x\rangle+x^{\top}Qx-\tr Q$,
\[
\log\E e^{itz}=it\,(m-\tr Q)-\tfrac12\log\det(I-2itQ)-\tfrac{t^2}2\,f^{\top}(I-2itQ)^{-1}f ,
\]
\begin{align*}
\kap_r(z)&=2^{r-1}(r-1)!\,\Tr Q^r+r!\,2^{r-3}f^{\top}Q^{r-2}f\\
&=2^{r-1}(r-1)!\,\Tr(D_cG)^r+r!\,2^{r-5}\,\omega^{\top}G(D_cG)^{r-2}G\omega :
\end{align*}
closed walks on the hidden codegree graph, and open walks whose two ends are the linear token.
(ii) Exactly (all chaoses), $\kap_r(z_{2,c})=\sum_{a_1,\dots,a_r}\prod_j\omega_{a_j}|w_{a_j}|\;\kap\big(\relu(s_{a_1}),\dots,\relu(s_{a_r})\big)$, and for
distinct indices the joint cumulant of the ridge functions is the sum over connected multigraphs on $\{a_1,\dots,a_r\}$
with weight $\rho^m/m!$ on an edge of multiplicity $m$ and weight $c_{\deg(v)}$ on a vertex of degree $\deg(v)$ (the diagram
formula). For $r=3$ the leading graphs are the three paths, centred at a hub of degree two (a kink, $c_2=\varphi_0$) with
two gated arms ($c_1=\tfrac12$), and the triangle ($c_2^3$). Repeated indices (two or three tokens on the same hidden
neuron: contacts) are given by the exact single-neuron and pair laws, not by the quadratic truncation.
\end{theorem}
\begin{proof}
(i) Diagonalise $Q$; each mode is a shifted noncentral $\chi^2$, and completing the square gives the displayed
log-characteristic function. Expand $\log\det(I-2itQ)=-\sum_r(2it)^r\Tr Q^r/r$ and $(I-2itQ)^{-1}=\sum_s(2it)^sQ^s$, and read off $r!$
times the coefficient of $(it)^r$. Substitute $Q=W_1^{\top}D_cW_1$ and $f=\tfrac12W_1^{\top}\omega$. (ii) Mehler's formula for
Gaussian vectors with unit variances: $\E\prod_jg_j(s_j)=\sum\prod_{v}\E[g_v^{(\deg v)}]\prod_{e}\rho_e^{m_e}/m_e!$ over multigraphs;
cumulants keep the connected ones.
\end{proof}
\begin{center}\small
\begin{tabular}{lccccc}\toprule
width & $\kap_3(z_{2,c})$ (MC) & self & contacts & distinct neurons & leading diagrams\\\midrule
$24$ & $-0.521$ & $-0.465$ & $-0.115$ & $0.059$ & $0.064$\\
$96$ & $0.210$ & $0.078$ & $0.121$ & $0.0115$ & $0.0133$\\\bottomrule
\end{tabular}
\end{center}
(Self: all three indices on one neuron; contacts: two of three on one neuron, by the exact pair law; distinct: Monte
Carlo minus the two exact parts; leading diagrams: hub paths and triangles.) The distinct-neuron part is the walk sum, to the accuracy of its leading diagrams; at these widths the contacts carry
most of a single pre-activation's third cumulant. The quadratic truncation of (i) applied to the whole of $z_{2,c}$ gives
$0.070$ against the true $-0.521$ ($n=24$): the walk picture is right for exchange between distinct neurons and wrong for
contacts, which must be computed exactly.

\subsection{How many tokens: depth heats the tower}
\begin{proposition}[Mean token number]\label{prop:tokens}
For $F\in L^2(\gamma)$ the mean token number of its centred part is $\bar N(F)=\E|\nabla F|^2/\Var F$ (since
$\E|\nabla F|^2=\sum_kk\,k!\|f_k\|^2$ and $\Var F=\sum_{k\ge1}k!\|f_k\|^2$). In the infinite-width limit the pre-activations of layer $l$
have
\[
\bar N_l=\frac1{1-\cos\theta_{l-1}},\qquad \theta_0=\tfrac\pi2,\quad \cos\theta_{k+1}=J(\theta_k)/\pi ,
\qquad \bar N_l\simeq\frac{2l^2}{9\pi^2}\ (l\to\infty),
\]
the pair temperature of Proposition~\ref{prop:heattrace}.
\end{proposition}
\begin{proof}
For a one-homogeneous random field with kernel $K(x,y)=\tfrac{|x||y|}nk(\hat x\cdot\hat y)$ differentiating at $y=x$, $|x|^2=n$, gives
$\E|\nabla_xz|^2=\sum_i\partial_{x_i}\partial_{y_i}K|_{y=x}=\tfrac1n\big(k(1)+(n-1)k'(1)\big)$ while $\E z^2=k(1)$; the normalised dual kernel of \textsc{ReLU} has derivative $1$ at $t=1$ at every layer (criticality), so
$\E|\nabla z|^2=\E z^2(1+O(1/n))$. The squared mean $m^2$ is the kernel at two independent inputs, whose cosine tends to $0$, so
$m^2/\E z^2\to\cos\theta_{l-1}$ with the angle recursion of Theorem~\ref{thm:folding} started at $\pi/2$. Then
$\bar N=\E z^2/(\E z^2-m^2)$; the asymptotics follow from $\beta_l\simeq9\pi^2/(2l^2)$.
\end{proof}
\begin{center}\small
\begin{tabular}{lccccccc}\toprule
layer & $1$ & $2$ & $3$ & $5$ & $9$ & $13$ & $16$\\\midrule
infinite width & $1.00$ & $1.47$ & $1.98$ & $3.13$ & $6.04$ & $9.72$ & $12.99$\\
width $256$ (Monte Carlo) & $1.00$ & $1.47$ & $2.01$ & $3.12$ & $6.26$ & $9.46$ & $10.56$\\\bottomrule
\end{tabular}
\end{center}
So the token tower does not truncate: the last pre-activation of a sixteen-layer network carries eleven to thirteen
tokens on average, far from the two of a quadratic form (and far below the $2^{l-1}$ of naive degree doubling).
Its marginal law is nevertheless close to Gaussian (skewness about $0.1$ on network $0$). The two facts are reconciled by the
fourth-moment theorem of Nualart and Peccati \cite{NP}: within a chaos, Gaussianity is controlled by the contractions $f\otimes_rf$
($1\le r<k$), the overlaps of a kernel's token configurations with themselves, i.e.\ token codegrees, and not by the token
number. The Gaussian closure works because these contractions are small at He initialisation, and it fails exactly
where they are not: at contacts, and along the collective mode.

\subsection{Two tiers: the input law is the environment}
\begin{theorem}[One-token plus two-token]\label{thm:tiers}
For a one-homogeneous $F$ and every Gaussian input law $N(m,\Sigma)$,
\[
\E_{N(m,\Sigma)}F=\big\langle m,\ \E\nabla F\big\rangle+\big\langle\Sigma,\ \E\nabla^2F\big\rangle ,
\]
with the distributional Hessian (walls). Equivalently, the exact mean $u(m,\Sigma)$ satisfies Euler's relation and the heat
equation $\partial_{\Sigma_{ij}}u=\tfrac12\partial_{m_i}\partial_{m_j}u$; in token language, adding a token pair at $(i,j)$ to the input
state has the same effect as two successive displacements along $i$ and $j$.
\end{theorem}
\begin{proof}
$\E F=\E[x\cdot\nabla F]=\langle m,\E\nabla F\rangle+\E[(x-m)\cdot\nabla F]$ and Stein's identity for $N(m,\Sigma)$,
$\E[(x-m)g]=\Sigma\,\E\nabla g$, with $g=\nabla F$. The heat equation is $\partial_\Sigma\E F(m+\Sigma^{1/2}Z)=\tfrac12\E\nabla^2F$.
\end{proof}

The upper tier is the input law; it sets the environment of the lower tier, namely the kink densities and the
gate-pair kernels of every layer. A moment chain $e(m,\Sigma)$ obeys Euler's relation (scale covariance) but not the heat
equation: its pair response and its double-displacement response disagree, and that disagreement is the heat defect of
stage 9. Along a localization (Eldan's balls of radius $r$ about a moving centre, Theorem~\ref{thm:uniform}) the environment
cools. Gates freeze, the contact weights $\Phi(1-\Phi)$ vanish away from walls, the two-token term transfers into the
one-token term, and at $r=0$ the walk is free and every chain is exact. The Euler--Stein correction integrates the
chain's mismatch over this cooling, and its constant $1/(d+2)$ records the token order of the mismatch at the kink.

\subsection{The two-token walk and its environment}
Pathwise, the gradient Gram $\Gamma_l=J_lJ_l^{\top}$ ($J_l=\partial z_l/\partial x$) obeys $\Gamma_l=W_lD_{l-1}\Gamma_{l-1}D_{l-1}W_l^{\top}$: two tokens
walk up the layered graph, each picking up weights and gates. In a factorised environment its expectation is the
completely positive map $\Gamma\mapsto W(K\circ\Gamma)W^{\top}$, where the gate-pair kernel $K(a,b)=\mathbb P(z_a>0,z_b>0)$ is a positive definite
function on the pair groupoid. The tetrachoric (Mehler) series splits it into three interactions:
\[
K=\underbrace{\Phi\Phi^{\top}}_{\text{independent walkers}}+\underbrace{\diag(\Phi-\Phi^2)}_{\text{contact}}+\underbrace{\sum_{k\ge1}\frac{1}{k!}\big(\rho^{\circ k}\circ\psi_k\psi_k^{\top}\big)_{\rm off}}_{\text{exchange of $k$ quanta}},
\qquad\psi_k=\varphi(\alpha)\He_{k-1}(\alpha).
\]
The first term is a congruence (the two tokens move independently with mean gates). The second acts only when both
tokens sit on one neuron, and it is diagonal, so it is sparse in the two-token graph. The third is dense, with
off-diagonal entries of order $n^{-k/2}$. The mean reads the unit space of the final quadratic form (its trace).

\begin{proposition}[Independent walkers are wrong at order one, by exactly the contact term]\label{prop:contactterm}
For two layers, let $u^{\rm walk}_c=\Phi_cm_c+f_c(0)\,\E|\nabla z_{2,c}|^2$ be the kink current with independent walkers, and
$u^{\rm GC}_c=\Phi_cm_c+f_c(0)\sigma_c^2$ the Gaussian closure (same exact $m_c,\sigma_c$, $\alpha_c=m_c/\sigma_c$). Then
\[
u^{\rm walk}_c-u^{\rm GC}_c=f_c(0)\big(\E|\nabla z_{2,c}|^2-\sigma_c^2\big)
=f_c(0)\sum_aW_2(c,a)\,\mu_a\big(m_c-\E[z_{2,c}\mid z_{1,a}=0]\big),
\]
whose contact part ($b=a$ in $z_{2,c}=\sum_bW_2(c,b)h_{1,b}$) is $f_c(0)\sum_aW_2(c,a)^2\mu_a^2=O(1)$, while the exchange part
($b\neq a$) is second order in $\rho_{ab}$.
\end{proposition}
\begin{proof}
Stein's identity gives $\E z^2=\E|\nabla z|^2+\E[z\Delta z]$ for the one-homogeneous $z=z_{2,c}$, and
$\Delta z=\sum_aW_2(c,a)\delta(z_{1,a})|w_a|^2$ with $f_{1,a}(0)|w_a|^2=\mu_a$; hence $\E|\nabla z|^2-\sigma^2=m^2-\sum_aW_2(c,a)\mu_a\E[z\mid z_{1,a}=0]$, and
$m=\sum_aW_2(c,a)\mu_a$. Conditioning on $z_{1,a}=0$ sets $h_{1,a}=0$ (the contact) and shifts the other neurons through their
correlations with $z_{1,a}$, which at zero mean changes their means at order $\rho^2$.
\end{proof}
Numerically (width $256$, $2^{20}$ inputs): the Gaussian closure has rms error $1.7\cdot10^{-3}$, independent walkers
$1.9\cdot10^{-1}$; $\E|\nabla z|^2-\sigma^2$ and the contact sum agree in mean ($0.639$ and $0.639$, correlation $0.999$); removing the contact term
brings the walkers to within $7.8\cdot10^{-4}$ of the Gaussian closure. So the Gaussian closure is the two-token walk with
exact contacts. Re-projecting each layer onto a fresh one-token (Gaussian) vector with the same covariance is precisely
what handles self-pinning, and the third-cumulant chain adds first-order exchange.

\subsection{What sparsifies, and what does not}
\begin{enumerate}[leftmargin=*]
\item \textbf{Two-token objects close.} Every sum
$\sum_{l'}T_{l\leftarrow l'}X_{l'}T_{l\leftarrow l'}^{\top}$ accumulates into one $n\times n$ recursion ($n^3$ per layer): the covariance, the gradient Gram, the Gaussian closure, and any hub sum whose hub coefficient is known
when the hub is born.
\item \textbf{Layer two is cheap.} The open walk of length two (the star of Theorem~\ref{thm:walks}) is
\[\textstyle\sum_aD_c(a)\,(W_2G)(c,a)^2 ;\] its hub coefficient is $W_2$ itself, known at birth, so all $n$ of them cost one product
$W_2G$ and a Hadamard square.
\item \textbf{From layer three on the hub coefficient is a walk.} A hub born at layer $l'$ and read at layer $l$ needs
three tokens that meet at the hub, each transported from $l'$ to $l$ (two arms and the reader). Nothing accumulates, and
this is the third-cumulant memory. The trace over input slots (the operation the Gaussian input performs) contracts
input vertices, never hidden hubs, so no operation on the input removes it.
\end{enumerate}
\begin{center}\small
\begin{tabular}{llc}\toprule
token-walk level & estimator & network $0$, MSE\\\midrule
independent walkers & factorised kink current & $O(1)$ (two layers)\\
$+$ exact contacts & Gaussian closure & $4.1\cdot10^{-6}$\\
$+$ first-order exchange & third-cumulant chain (radial $\kap_4$ off/on) & $5.6\cdot10^{-7}$ / $3.5\cdot10^{-8}$\\
$+$ cooling, $d=2$ & Euler--Stein merge, $a\approx\tfrac14$ & $6.6\cdot10^{-9}$\\\bottomrule
\end{tabular}
\end{center}

\section{What the placement explains, what it predicts, and what is open}
\label{sec:outlook}
\subsection{What the placement explains}
\begin{enumerate}[leftmargin=*]
\item \textbf{Why the mean is hard and the densities are easy.} The mean is a sum over walls of a density at the kink
times a conditional transport on the wall (Theorem~\ref{thm:kink}). Densities are concentration-function
coefficients and the Gaussian surrogate gets them right; the conditional transports are rare-event-relative objects
(conditioning on a measure-zero wall), and every chain's error lives there.
\item \textbf{Why the third-cumulant memory has an identity leg.} It is the label of the kink at which the cumulant was
born (Theorem~\ref{thm:birth}); the memory is the record of all pinnings, transported.
\item \textbf{Why it cannot be compressed.} Kinks have small codegree at He initialisation, so their atoms are orthogonal;
the frame bound (Theorem~\ref{thm:frame}) makes the memory exactly as compressible as its transport leg, and the
complementarity of the neuron algebra with the modular frame of the layer state (Theorem~\ref{thm:complement}) makes
state-respecting (spectral) compressions blind to the hub diagonal up to $2\sqrt{2/n}$. The same small codegree that
makes Sahasrabudhe's nibble concentrate makes our compression fail.
\item \textbf{Why ensemble surrogates fail.} The memory's readout is odd in the last weight matrix
(Proposition~\ref{prop:odd}); free and NNGP descriptions are the symmetric phase.
\item \textbf{Where the Euler--Stein constants come from.} The correction averages the defect over Gaussian balls of
uniform radius (Theorem~\ref{thm:uniform}); a defect of Hermite contact order $d$ decays like $r^{d+1}$, so
$a=1/(d+2)$, and $d$ is the order of the concentration coefficient the estimator gets wrong
(Theorem~\ref{thm:order}, Proposition~\ref{prop:conc}); the three measured constants fit $d=0,0,2$.
\item \textbf{Why depth behaves algebraically.} Lengths are preserved and chords contract like $3\pi/l$
(Theorem~\ref{thm:folding}) because the activation's heat trace is critical at linear order and has a
$\beta^{3/2}$ kink term (Proposition~\ref{prop:heattrace}).
\end{enumerate}

\subsection{What it predicts (falsifiable, theoretical)}
\begin{enumerate}[leftmargin=*]
\item The defect of the exact first-order chain along Eldan's path decays like $r^3=(1+\tau)^{-3/2}$ (order $d=2$),
not like $(1+\tau)^{-1}$; a chain that also corrects the curvature coefficient $f''(0)$ at the kinks has a defect of order
$d\ge4$ and merge constant at most $1/6$.
\item On the official networks, the number of walls crossed by a great-circle arc grows linearly in depth (about
$n\theta_0/\pi$ per layer) while the number of gates that differ between its endpoints grows like $3n\log L$.
\item The best hub-index compression of any old source has the error of the low-rank truncation of its transport leg
$T_{l\leftarrow l'}C$ (Theorem~\ref{thm:frame} with a near-isotropic Hadamard leg); a tier cannot do better than that number,
which can be computed without running any chain.
\item The complementarity constant of the layer state at depth is $\approx2\sqrt{2/n}$ up to the collective mode, so every
state-preserving coarse-graining of a layer onto spectral data retains at most a fraction $\approx8/n$ of the centred
energy of any neuron-diagonal observable (a gate pattern, a hub diagonal).
\end{enumerate}

\subsection{Open problems}
\begin{enumerate}[leftmargin=*]
\item \textbf{Contact order at depth.} Prove Conjecture~\ref{conj:order}: show that the defect of a chain whose first
neglected term is of Hermite order $d$ at every wall has the profile $r^{d+1}$ along the localization, despite the
nonlinear dependence of the downstream pre-activation means on the localization centre.
\item \textbf{A contact-resolved estimator.} The kink-current formula asks for the on-wall transports
$\E[\partial_{h_l}h_L\,e_a\mid z_{l,a}=0]$ for all walls. To first order these are pinning responses
(Lemma~\ref{lem:pinning}); their sum over walls is one directional derivative of the downstream chain per output,
which forward mode computes with an identity leg (the memory) and reverse mode with one adjoint per output. Either way
the cost is that of the memory. Is there an order-two quantity (a pinning response of pinning responses) whose sum
over walls collapses, as the radial fourth-cumulant scalar does for the fourth cumulant?
\item \textbf{The induced Dirac ensemble.} Compute the spectral dimension and spectral variance \cite{BDG} of the
selected Dirac operators $P_xDP_x$ and of the transports $T_{l\leftarrow l'}$ by free multiplicative convolution, and
compare with the measured stable ranks $164,92,65,47,29,23$ of the cocycle at ages $1,2,3,4,6,8$; decide whether the
emergence of the collective eigenvalue of the layer covariance (from $0.9\%$ to $37$--$48\%$ of the energy) is a
Baik--Ben Arous--P\'ech\'e transition in depth, the analogue of the one-cut to two-cut transition of \cite{DG}.
\item \textbf{Bootstrap with kink generators.} Is the island of the moment problem at low degree narrow at large $n$
because of a factorisation analogous to large-$N$ factorisation \cite{KP}? If so, positivity would certify chains
cheaply.
\item \textbf{Klartag in reverse.} Klartag's ellipsoid stops when about $d^2/2$ contacts hold its $d(d+1)/2$ shape
parameters. The marginal value of a source age decays by a factor $0.75$ per age and about eleven ages matter,
roughly $4n$ effective contacts held in memory. Is there a counting argument that fixes this number from the dimension
of the third-cumulant readouts that the mean actually uses?
\end{enumerate}

\section*{Verification}
\label{sec:verify}
\texttt{scripts/verify\_stage10.py} compares the two sides of each identity numerically; it runs no estimator and
no network of the competition. Output in \texttt{notes/stage10/verify\_stage10.txt}.
\begin{center}\small
\begin{tabular}{>{\raggedright\arraybackslash}p{4.4cm}>{\raggedright\arraybackslash}p{6.0cm}>{\raggedright\arraybackslash}p{4.4cm}}\toprule
statement & check & result\\\midrule
Theorem~\ref{thm:kink} (kink current) & $\E h_L$ against the wall sum with a $\pm0.05$ window, $3\cdot10^6$ inputs, $L=3$ & $n=16$: agreement to noise ($1.5\%$); $n=6$: $25\%$ off, dead tiles carry $16\%$ \\
Theorem~\ref{thm:crofton} (Crofton) & crossings of $2\cdot10^4$ random walls by a random curve in $\R^{50}$ & $31.28$ against $\Lambda/\pi=31.39$\\
Theorem~\ref{thm:folding}, Proposition~\ref{prop:heattrace} & iterate the angle map & $\theta_l l/3\pi=0.90,0.985,0.992$ at $l=10^2,10^3,2\cdot10^3$ (log corrections)\\
Proposition~\ref{prop:gain} (unit gain) & tangent and norm gains at $\sigma_w^2=2$, $n=400$ & $1.009$, $1.002$\\
Proposition~\ref{prop:mehler} (Mehler, Hermite coefficients of $\relu$) & closed form against the Hermite series & equal to $10^{-8}$\\
Theorem~\ref{thm:order} (contact order) & stage-9 $\tau$-integral for $d=0,2,4$ & $0.497,0.246,0.164$ against $\tfrac12,\tfrac14,\tfrac16$ (truncated tail)\\
Theorem~\ref{thm:birth} (births are contacts) & $\kap_3/\varepsilon^2$ by Sobol quasi-Monte Carlo, $2^{23}$ points & ratio to the hub formula $1.064,1.030,1.015,1.009$ at $\varepsilon=0.16,\dots,0.02$\\
Theorem~\ref{thm:complement} (complementarity) & $\|E_BE_A-J\|$ against $\|C\|$ for $S=WW^{\top}$ & equal; $0.333$ / $0.173$ against $2\sqrt{2/n}=0.354$ / $0.177$ ($n=64,256$)\\
Theorem~\ref{thm:frame} (frame bound) & $2000$ random rank-$10$ compressions & bound holds\\\bottomrule
\end{tabular}
\end{center}
The asymptotic $3\pi/l$ is reached slowly ($O(\log l/l)$ corrections), so at the depth $L=16$ of the competition the exact
recursion, not its asymptotic form, is the statement to use.

\begin{thebibliography}{99}\small
\bibitem{stage9} Stage 9 note: \emph{the mean of a random ReLU network as a central localization} (this repository, \texttt{notes/stage9}).
\bibitem{stage8} Stage 8 note: \emph{Localization in the commutant: modular tilts, Kikuchi--Toeplitz quantization and transverse measures}.
\bibitem{survey} J.~Sahasrabudhe, Probabilistic combinatorics at exponentially small scales, arXiv:2512.15077.
\bibitem{klartag} B.~Klartag, Lattice packing of spheres in high dimensions using a stochastically evolving ellipsoid, arXiv:2504.05042.
\bibitem{DG} M.~D'Arcangelo, S.~Gnutzmann, Symmetry breaking and phase transitions in random non-commutative geometries and related random-matrix ensembles, arXiv:2601.14141.
\bibitem{KP} M.~Khalkhali, N.~Pagliaroli, Bootstrapping noncommutative geometry with Dirac ensembles, arXiv:2512.08694.
\bibitem{BDG} J.~W.~Barrett, P.~Druce, L.~Glaser, Spectral estimators for finite non-commutative geometries, arXiv:1902.03590.
\bibitem{BG} J.~W.~Barrett, L.~Glaser, Monte Carlo simulations of random non-commutative geometries, J.~Phys.~A 49 (2016) 245001.
\bibitem{WLZ} R.~Wang, L.~C.~Lau, H.~Zhou, Derandomizing matrix concentration inequalities from free probability, arXiv:2601.08111.
\bibitem{SST} A.~Sah, J.~Sahasrabudhe, M.~Sawhney, On the Spielman--Teng conjecture, arXiv:2405.20308.
\bibitem{CS} Y.~Cho, L.~K.~Saul, Kernel methods for deep learning, NeurIPS 2009.
\bibitem{PLRSG} B.~Poole, S.~Lahiri, M.~Raghu, J.~Sohl-Dickstein, S.~Ganguli, Exponential expressivity in deep neural networks through transient chaos, NeurIPS 2016.
\bibitem{HDR} S.~Hayou, A.~Doucet, J.~Rousseau, On the impact of the activation function on deep neural networks training, ICML 2019.
\bibitem{HR} B.~Hanin, D.~Rolnick, Complexity of linear regions in deep networks, ICML 2019.
\bibitem{Scarf} H.~Scarf, A min-max solution of an inventory problem, in \emph{Studies in the Mathematical Theory of Inventory and Production}, Stanford 1958.
\bibitem{NP} D.~Nualart, G.~Peccati, Central limit theorems for sequences of multiple stochastic integrals, Ann.~Probab.~33 (2005) 177--193; I.~Nourdin, G.~Peccati, \emph{Normal Approximations with Malliavin Calculus}, Cambridge 2012.
\bibitem{Las} J.-B.~Lasserre, Global optimization with polynomials and the problem of moments, SIAM J.~Optim.~11 (2001) 796--817.
\bibitem{Tak} M.~Takesaki, Conditional expectations in von Neumann algebras, J.~Funct.~Anal.~9 (1972) 306--321.
\bibitem{Rief} M.~A.~Rieffel, Metrics on state spaces, Doc.~Math.~4 (1999) 559--600.
\bibitem{Connes} A.~Connes, \emph{Noncommutative Geometry}, Academic Press 1994.
\end{thebibliography}
\end{document}
